\ifdefined\HAFOFractionalOnly\else
\section{Sistemas estudiados}
\label{sec:modelos}

\subsection{Sistema de Chua entero de referencia}

El sistema entero utilizado para recuperar el método de
Leonov--Kuznetsov es
\begin{equation}
\begin{aligned}
 \dot x &=\alpha\bigl(y-x-f(x)\bigr),\\
 \dot y &=x-y+z,\\
 \dot z &=-\beta y-\gamma z,\\
 f(x)&=m_1x+(m_0-m_1)\operatorname{sat}(x),
\end{aligned}
\label{eq:chua_entero}
\end{equation}
donde
\begin{equation}
\operatorname{sat}(x)=
\begin{cases}
-1,&x<-1,\\
x,&|x|\leq1,\\
1,&x>1.
\end{cases}
\label{eq:saturacion}
\end{equation}
El caso de referencia usa
\begin{equation}
(\alpha,\beta,\gamma,m_0,m_1)
=(8.4562,12.0732,0.0052,-0.1768,-1.1468).
\label{eq:parametros_enteros}
\end{equation}

\subsection{Chua fraccionario no suave}

El primer sistema de estudio sustituye las tres derivadas ordinarias de
\cref{eq:chua_entero} por derivadas de Caputo de orden conmensurable:
\begin{equation}
\begin{aligned}
 {}^{\mathrm C}D_{t_0}^{q}x
 &=\alpha\bigl(y-x-f_{\mathrm{ns}}(x)\bigr),\\
 {}^{\mathrm C}D_{t_0}^{q}y
 &=x-y+z,\\
 {}^{\mathrm C}D_{t_0}^{q}z
 &=-\beta y-\gamma z,\\
 f_{\mathrm{ns}}(x)
 &=m_1x+(m_0-m_1)\operatorname{sat}(x),
\qquad 0<q\leq1.
\end{aligned}
\label{eq:chua_no_suave}
\end{equation}
Los parámetros nominales y el orden son
\begin{equation}
\begin{aligned}
q&=0.9998, &
(\alpha,\beta,\gamma)&=(8.4562,12.0732,0.0052),\\
(m_0,m_1)&=(-0.1768,-1.1468).
\end{aligned}
\label{eq:parametros_no_suave}
\end{equation}

\subsection{Chua fraccionario con arctangente}

El segundo sistema de estudio es
\begin{equation}
\begin{aligned}
 {}^{\mathrm C}D_{t_0}^{q}x
 &=\alpha\bigl(y-x-h_{\mathrm a}(x)\bigr),\\
 {}^{\mathrm C}D_{t_0}^{q}y
 &=x-y+z,\\
 {}^{\mathrm C}D_{t_0}^{q}z
 &=-\beta y-\gamma z,\\
 h_{\mathrm a}(x)&=a_1x+a_2\arctan(\rho x).
\end{aligned}
\label{eq:chua_arctan}
\end{equation}
La reconstrucción de Wu et al. usa
\begin{equation}
\begin{aligned}
q&=0.99,&
(\alpha,\beta,\gamma)&=(8.4562,12.0732,0.0052),\\
a_1&=0.4,&a_2&=-1.5585,&\rho&=1.
\end{aligned}
\label{eq:parametros_wu}
\end{equation}
El segundo caso de esta familia usa
\begin{equation}
\begin{aligned}
q&=0.9999,&
\alpha&=21.8493569066,&
\beta&=19.0818408409,\\
\gamma&=0.0073780120,&
a_1&=0.0422897934,&
a_2&=-3.3367815123,\\
\rho&=1.7984259333.
\end{aligned}
\label{eq:parametros_c590}
\end{equation}
Los dos conjuntos de parámetros pertenecen a la misma familia
\cref{eq:chua_arctan}, pero corresponden a experimentos independientes.

Los equilibrios, los jacobianos y los polinomios característicos determinan
los datos espectrales necesarios para aplicar el criterio de estabilidad
fraccionaria.

\begingroup
\def\HAFOderivationpart{3}
% Este archivo es una biblioteca de derivaciones. Cada bloque se inserta en el
% lugar del cuerpo donde se usa, mediante \HAFOderivationpart. No se imprime
% como apéndice independiente.
\ifnum\HAFOderivationpart=1\relax
\subsection{Derivación completa de la forma de Lur'e}
\label{app:verificacion_lure}

Sea \(h(x)=\ell x+\psi(x)\). La primera ecuación de los dos modelos es
\[
{}^{\mathrm C}D_{t_0}^{q}x
=\alpha(y-x-h(x)).
\]
Al distribuir \(\alpha\),
\begin{align}
\alpha(y-x-h(x))
&=\alpha y-\alpha x-\alpha\ell x-\alpha\psi(x)\nonumber\\
&=-\alpha(1+\ell)x+\alpha y-\alpha\psi(x).
\label{eq:expansion_primera_fila}
\end{align}
Por otra parte,
\[
PX=
\begin{pmatrix}
-\alpha(1+\ell)x+\alpha y\\
x-y+z\\
-\beta y-\gamma z
\end{pmatrix},
\qquad
b\psi(r^{\T}X)=
\begin{pmatrix}
-\alpha\psi(x)\\0\\0
\end{pmatrix}.
\]
La suma reproduce \cref{eq:expansion_primera_fila} y las otras dos
ecuaciones. Para la saturación se toma
\(\ell=m_1\) y
\(\psi=(m_0-m_1)\operatorname{sat}\); para la arctangente,
\(\ell=a_1\) y \(\psi=a_2\arctan(\rho\,\cdot)\).

\fi
\ifnum\HAFOderivationpart=2\relax
\subsection{Derivación completa del determinante y la transferencia}
\label{app:transferencia_explicita}

Con \(\lambda=s^q\),
\begin{equation}
\lambda I-P=
\begin{pmatrix}
\lambda+\alpha(1+\ell)&-\alpha&0\\
-1&\lambda+1&-1\\
0&\beta&\lambda+\gamma
\end{pmatrix}.
\label{eq:matriz_resolvente}
\end{equation}
El menor de la posición \((1,1)\) es
\begin{align}
Q(\lambda)
&=
\det
\begin{pmatrix}
\lambda+1&-1\\
\beta&\lambda+\gamma
\end{pmatrix}\nonumber\\
&=(\lambda+1)(\lambda+\gamma)+\beta.
\label{eq:menor_Q}
\end{align}
Al desarrollar \(\det(\lambda I-P)\) por la primera fila,
\begin{align}
\Delta_\ell(\lambda)
&=\bigl[\lambda+\alpha(1+\ell)\bigr]Q(\lambda)
-(-\alpha)
\det
\begin{pmatrix}
-1&-1\\
0&\lambda+\gamma
\end{pmatrix}\nonumber\\
&=\bigl[\lambda+\alpha(1+\ell)\bigr]Q(\lambda)
-\alpha(\lambda+\gamma).
\label{eq:determinante_transferencia}
\end{align}

Para calcular la primera componente de
\((\lambda I-P)^{-1}b\), la regla de Cramer sustituye la primera columna de
\cref{eq:matriz_resolvente} por
\(b=(-\alpha,0,0)^{\T}\). El determinante del numerador es
\[
\det
\begin{pmatrix}
-\alpha&-\alpha&0\\
0&\lambda+1&-1\\
0&\beta&\lambda+\gamma
\end{pmatrix}
=-\alpha Q(\lambda).
\]
Como \(r^{\T}\) selecciona esa primera componente,
\begin{equation}
W_q(s)
=r^{\T}(s^qI-P)^{-1}b
=-\alpha\frac{Q(s^q)}{\Delta_\ell(s^q)}.
\label{eq:transferencia_cramer}
\end{equation}

La identidad
\[
P-s^qI=-(s^qI-P)
\]
implica
\[
r^{\T}(P-s^qI)^{-1}b=-W_q(s).
\]
Por ello, cambiar el orden de los términos dentro de la resolvente cambia
simultáneamente el signo de la transferencia y el signo del cierre. No cambia
la solución física si ambas modificaciones se realizan de manera coherente.

\fi
\ifnum\HAFOderivationpart=3\relax
\subsection{Derivación completa de los equilibrios}
\label{app:equilibrios}

En un equilibrio \(e=(x_e,y_e,z_e)^{\T}\), las dos últimas ecuaciones de
ambos sistemas dan
\[
0=x_e-y_e+z_e,\qquad
0=-\beta y_e-\gamma z_e.
\]
De la primera,
\[
z_e=y_e-x_e.
\]
Al sustituirla en la segunda,
\begin{align*}
0
&=-\beta y_e-\gamma(y_e-x_e)\\
&=-(\beta+\gamma)y_e+\gamma x_e.
\end{align*}
Si \(\beta+\gamma\neq0\),
\begin{equation}
y_e=\frac{\gamma}{\beta+\gamma}x_e,\qquad
z_e=-\frac{\beta}{\beta+\gamma}x_e.
\label{eq:equilibrios_yz}
\end{equation}
La primera ecuación de equilibrio exige
\[
0=y_e-x_e-h(x_e).
\]
Mediante \cref{eq:equilibrios_yz},
\begin{equation}
h(x_e)+\frac{\beta}{\beta+\gamma}x_e=0.
\label{eq:equilibrio_escalar_comun}
\end{equation}
Esta única ecuación escalar determina todos los equilibrios; las otras dos
coordenadas se recuperan después con \cref{eq:equilibrios_yz}.

\subsubsection{Equilibrios del Chua no suave}

En la región \(|x_e|<1\), \(h(x_e)=m_0x_e\). Entonces
\begin{equation}
\left(m_0+\frac{\beta}{\beta+\gamma}\right)x_e=0.
\label{eq:equilibrio_interior_ns}
\end{equation}
Cuando el coeficiente no es cero, la única raíz interior es \(x_e=0\).

Para \(x_e>1\),
\[
h(x_e)=m_1x_e+(m_0-m_1),
\]
y \cref{eq:equilibrio_escalar_comun} queda
\[
\left(m_1+\frac{\beta}{\beta+\gamma}\right)x_e+(m_0-m_1)=0.
\]
Por tanto,
\begin{equation}
x_\star=
\frac{m_1-m_0}
{m_1+\dfrac{\beta}{\beta+\gamma}}.
\label{eq:raiz_exterior_ns}
\end{equation}
Si \(x_\star>1\), existe el equilibrio positivo. Como la no linealidad es
impar, existe también la raíz \(-x_\star<-1\). Los tres equilibrios son
\begin{equation}
E_0=(0,0,0),\qquad
E_\pm=
\left(
\pm x_\star,\,
\pm\frac{\gamma}{\beta+\gamma}x_\star,\,
\mp\frac{\beta}{\beta+\gamma}x_\star
\right).
\label{eq:equilibrios_ns_simbolicos}
\end{equation}
La condición regional \(x_\star>1\) debe comprobarse; no basta con resolver la
ecuación lineal exterior.

\subsubsection{Equilibrios del Chua con arctangente}

Al usar
\(h(x)=a_1x+a_2\arctan(\rho x)\) en
\cref{eq:equilibrio_escalar_comun}, se obtiene
\begin{equation}
\left(a_1+\frac{\beta}{\beta+\gamma}\right)x_e
+a_2\arctan(\rho x_e)=0.
\label{eq:equilibrio_escalar_arctan}
\end{equation}
La expresión es impar, de modo que \(x_e=0\) siempre es raíz y toda raíz
positiva no nula tiene una compañera negativa. Las raíces no nulas se
determinan numéricamente y se sustituyen en \cref{eq:equilibrios_yz}. Esta
ecuación se evalúa con \(\rho=1\) para el caso bibliográfico y con el valor de
\(\rho\) de \cref{eq:parametros_c590} para el caso \texttt{c590}; no deben
intercambiarse las dos parametrizaciones.

\fi
\ifnum\HAFOderivationpart=4\relax
\subsection{Derivación completa de los jacobianos regionales}
\label{app:jacobianos}

Para la forma de Lur'e,
\[
F(X)=PX+b\psi(r^{\T}X).
\]
La regla de la cadena produce
\begin{equation}
J_e=P+b\,\psi'(r^{\T}e)r^{\T}.
\label{eq:jacobiano_lure}
\end{equation}
Como \(r^{\T}e=x_e\), solamente cambia la entrada \((1,1)\).

En el sistema no suave,
\[
\psi'_{\mathrm{ns}}(x)=
\begin{cases}
m_0-m_1,&|x|<1,\\
0,&|x|>1.
\end{cases}
\]
Por ello,
\begin{equation}
J_{\mathrm{in}}=
\begin{pmatrix}
-\alpha(1+m_0)&\alpha&0\\
1&-1&1\\
0&-\beta&-\gamma
\end{pmatrix},
\qquad
J_{\mathrm{out}}=
\begin{pmatrix}
-\alpha(1+m_1)&\alpha&0\\
1&-1&1\\
0&-\beta&-\gamma
\end{pmatrix}.
\label{eq:jacobianos_ns}
\end{equation}
En \(x=\pm1\) la derivada clásica no existe. Los equilibrios usados en este
reporte no se encuentran sobre esas superficies, por lo que se emplea el
jacobiano de la región que contiene a cada uno.

Para la arctangente,
\[
\psi'_{\mathrm a}(x)
=\frac{a_2\rho}{1+\rho^2x^2},
\]
y
\begin{equation}
J_e=
\begin{pmatrix}
\displaystyle
-\alpha\left(1+a_1+
\frac{a_2\rho}{1+\rho^2x_e^2}\right)
&\alpha&0\\[3mm]
1&-1&1\\
0&-\beta&-\gamma
\end{pmatrix}.
\label{eq:jacobiano_arctan}
\end{equation}
Los valores propios de
\cref{eq:jacobianos_ns,eq:jacobiano_arctan} se clasifican con
\cref{eq:matignon}; la estabilidad local no sustituye el sondeo de cuenca.

\fi
\ifnum\HAFOderivationpart=5\relax
\subsection{Derivación completa del polinomio característico}
\label{app:polinomio_caracteristico}

Sea \(g=h'(x_e)\) la pendiente total de la no linealidad en un equilibrio.
El jacobiano de cualquiera de las dos familias puede escribirse como
\begin{equation}
J(g)=
\begin{pmatrix}
-\alpha(1+g)&\alpha&0\\
1&-1&1\\
0&-\beta&-\gamma
\end{pmatrix}.
\label{eq:jacobiano_pendiente_total}
\end{equation}
Por tanto,
\[
\lambda I-J(g)=
\begin{pmatrix}
\lambda+\alpha(1+g)&-\alpha&0\\
-1&\lambda+1&-1\\
0&\beta&\lambda+\gamma
\end{pmatrix}.
\]
El desarrollo por la primera fila, igual al de
\cref{eq:determinante_transferencia}, produce
\begin{equation}
\chi_g(\lambda)
=\bigl[\lambda+\alpha(1+g)\bigr]
\bigl[(\lambda+1)(\lambda+\gamma)+\beta\bigr]
-\alpha(\lambda+\gamma).
\label{eq:polinomio_caracteristico_agrupado}
\end{equation}
Como
\[
(\lambda+1)(\lambda+\gamma)+\beta
=\lambda^2+(1+\gamma)\lambda+(\beta+\gamma),
\]
la multiplicación y agrupación por potencias de \(\lambda\) da
\begin{equation}
\begin{aligned}
\chi_g(\lambda)
={}&\lambda^3+
\bigl[1+\gamma+\alpha(1+g)\bigr]\lambda^2\\
&+\bigl[
\beta+\gamma+\alpha(1+g)(1+\gamma)-\alpha
\bigr]\lambda\\
&+\alpha\bigl[(1+g)(\beta+\gamma)-\gamma\bigr].
\end{aligned}
\label{eq:polinomio_caracteristico_expandido}
\end{equation}
En el sistema no suave se sustituye \(g=m_0\) en \(E_0\) y \(g=m_1\)
en \(E_\pm\). En el sistema con arctangente,
\[
g=a_1+\frac{a_2\rho}{1+\rho^2x_e^2}.
\]
Las raíces de \cref{eq:polinomio_caracteristico_expandido} son los valores
propios utilizados en \cref{eq:matignon}; así, la clasificación local se
obtiene a partir del mismo campo que define la integración.

\fi
\ifnum\HAFOderivationpart=6\relax
\subsection{Derivación completa para la saturación centrada}
\label{app:df_saturacion_centrada}

Sea
\[
\psi_{\mathrm{ns}}(\sigma)=\Delta\operatorname{sat}(\sigma),
\qquad \Delta=m_0-m_1.
\]
Si \(0<A\leq1\), no hay saturación y
\[
\psi_{\mathrm{ns}}(A\sin\theta)=\Delta A\sin\theta.
\]
Como \(\int_0^\pi\sin^2\theta\dd\theta=\pi/2\),
\[
N_{\mathrm{ns}}(A)
=\frac{2}{\pi A}\Delta A\frac{\pi}{2}
=\Delta.
\]

Para \(A>1\), defínase
\[
\theta_0=\arcsin\left(\frac1A\right),
\qquad 0<\theta_0<\frac{\pi}{2}.
\]
En \(0\leq\theta\leq\theta_0\) y
\(\pi-\theta_0\leq\theta\leq\pi\), la entrada permanece en la parte lineal;
en el intervalo central la saturación vale \(1\). Entonces
\begin{align}
N_{\mathrm{ns}}(A)
&=\frac{2\Delta}{\pi A}
\left[
2A\int_0^{\theta_0}\sin^2\theta\dd\theta
+\int_{\theta_0}^{\pi-\theta_0}\sin\theta\dd\theta
\right]\nonumber\\
&=\frac{2\Delta}{\pi A}
\left[
A\theta_0-\frac{A}{2}\sin(2\theta_0)
+2\cos\theta_0
\right]\nonumber\\
&=\frac{2\Delta}{\pi A}
\left[A\theta_0+\cos\theta_0\right].
\label{eq:df_sat_pasos}
\end{align}
En el último paso se usó
\[
\frac{A}{2}\sin(2\theta_0)
=A\sin\theta_0\cos\theta_0
=\cos\theta_0.
\]
Finalmente,
\[
\cos\theta_0
=\sqrt{1-\frac1{A^2}}
=\frac{\sqrt{A^2-1}}{A},
\]
y \cref{eq:df_sat_pasos} se convierte en
\[
N_{\mathrm{ns}}(A)
=\Delta\frac{2}{\pi}
\left[
\arcsin\left(\frac1A\right)
+\frac{\sqrt{A^2-1}}{A^2}
\right],
\]
que coincide con \cref{eq:df_saturacion}.

\fi
\ifnum\HAFOderivationpart=7\relax
\subsection{Derivación completa de la saturación sesgada}
\label{app:df_saturacion_sesgada}

Considérese
\[
u(\theta)=c+A\cos\theta,\qquad A>0.
\]
Para escribir una fórmula que incluya los casos parcialmente y completamente
saturados, se define
\[
[v]_{-1}^{1}=\max(-1,\min(1,v))
\]
y los ángulos
\begin{equation}
\theta_+
=\arccos\left[\frac{1-c}{A}\right]_{-1}^{1},
\qquad
\theta_-
=\arccos\left[\frac{-1-c}{A}\right]_{-1}^{1}.
\label{eq:angulos_saturacion_sesgada}
\end{equation}
Como \((1-c)/A>(-1-c)/A\) y \(\arccos\) es decreciente,
\(0\leq\theta_+\leq\theta_-\leq\pi\). En el intervalo
\([0,\pi]\),
\[
\operatorname{sat}(u(\theta))=
\begin{cases}
1,&0\leq\theta\leq\theta_+,\\
c+A\cos\theta,&\theta_+\leq\theta\leq\theta_-,\\
-1,&\theta_-\leq\theta\leq\pi.
\end{cases}
\]
La simetría respecto de \(2\pi\) permite integrar solamente en
\([0,\pi]\). Para \(\psi=\Delta\operatorname{sat}\),
\begin{align}
\Psi_0(c,A)
=\frac{\Delta}{\pi}\Big[
&\theta_+
+c(\theta_--\theta_+)
+A(\sin\theta_--\sin\theta_+)\nonumber\\
&-(\pi-\theta_-)
\Big].
\label{eq:psi0_saturacion_cerrada}
\end{align}

Para la componente fundamental se usan
\[
\int\cos\theta\dd\theta=\sin\theta,\qquad
\int\cos^2\theta\dd\theta
=\frac{\theta}{2}+\frac{\sin2\theta}{4}.
\]
La sustitución por intervalos da
\begin{align}
\Psi_1(c,A)
=\frac{2\Delta}{\pi}\Bigg[
&\sin\theta_++\sin\theta_-
+c(\sin\theta_--\sin\theta_+)\nonumber\\
&+\frac{A}{2}(\theta_--\theta_+)
+\frac{A}{4}
\bigl(\sin2\theta_--\sin2\theta_+\bigr)
\Bigg].
\label{eq:psi1_saturacion_cerrada}
\end{align}
Las operaciones de recorte en
\cref{eq:angulos_saturacion_sesgada} hacen que
\cref{eq:psi0_saturacion_cerrada,eq:psi1_saturacion_cerrada} sigan siendo
válidas cuando alguna de las tres regiones tiene longitud cero.

La ecuación de media puede reducirse a una igualdad escalar. De
\cref{eq:transferencia_cramer},
\begin{align}
W_q(0)
&=-\alpha
\frac{\beta+\gamma}
{\alpha(1+\ell)(\beta+\gamma)-\alpha\gamma}\nonumber\\
&=-\frac{\beta+\gamma}
{(1+\ell)(\beta+\gamma)-\gamma}.
\label{eq:ganancia_dc}
\end{align}
Por tanto,
\begin{equation}
c=W_q(0)\Psi_0(c,A),
\label{eq:cierre_dc_escalar}
\end{equation}
y el cierre armónico es
\[
1-W_q(\ii\omega)\frac{\Psi_1(c,A)}{A}=0.
\]
Estas tres ecuaciones reales ---una para la media, y las partes real e
imaginaria del cierre--- determinan \((c,A,\omega)\).

\fi
\ifnum\HAFOderivationpart=8\relax
\subsection{Derivación completa para la arctangente}
\label{app:df_arctan}

Sea
\[
\psi_{\mathrm a}(\sigma)=a_2\arctan(\rho\sigma),
\qquad
\xi=\rho A.
\]
Se define
\[
I(\xi)=
\int_0^\pi\arctan(\xi\sin\theta)\sin\theta\dd\theta.
\]
Como \(I(0)=0\), puede calcularse \(I\) integrando su derivada:
\begin{align}
I'(\xi)
&=\int_0^\pi
\frac{\sin^2\theta}{1+\xi^2\sin^2\theta}\dd\theta\nonumber\\
&=\frac1{\xi^2}
\left[
\pi-\int_0^\pi
\frac{\dd\theta}{1+\xi^2\sin^2\theta}
\right]\nonumber\\
&=\frac{\pi}{\xi^2}
\left(1-\frac1{\sqrt{1+\xi^2}}\right).
\label{eq:derivada_integral_arctan}
\end{align}
En la última igualdad se usó la integral estándar
\[
\int_0^\pi\frac{\dd\theta}{1+\xi^2\sin^2\theta}
=\frac{\pi}{\sqrt{1+\xi^2}}.
\]
Una primitiva que se anula en \(\xi=0\) es
\begin{equation}
I(\xi)
=\pi\frac{\sqrt{1+\xi^2}-1}{\xi}.
\label{eq:integral_arctan}
\end{equation}
Al sustituir \(\xi=\rho A\) en
\cref{eq:df_centrada},
\begin{align}
N_{\mathrm a}(A)
&=\frac{2a_2}{\pi A}I(\rho A)\nonumber\\
&=\frac{2a_2}{\rho A^2}
\left(\sqrt{1+\rho^2A^2}-1\right)\nonumber\\
&=\frac{2a_2\rho}
{1+\sqrt{1+\rho^2A^2}}.
\label{eq:df_arctan_derivada}
\end{align}

Si el cierre exige \(N_{\mathrm a}(A)=k\), la amplitud puede despejarse.
De la forma racionalizada,
\[
1+\sqrt{1+\rho^2A^2}=\frac{2a_2\rho}{k}.
\]
Al aislar la raíz, elevar al cuadrado y simplificar,
\begin{equation}
A^2
=\frac{4a_2(a_2\rho-k)}{k^2\rho}.
\label{eq:amplitud_arctan}
\end{equation}
La solución se acepta solamente si el lado derecho es positivo y la igualdad
previa a elevar al cuadrado tiene lado derecho no menor que \(2\). De manera
equivalente, para \(A>0\), \(k\) tiene el mismo signo que \(a_2\) y se
encuentra estrictamente entre \(0\) y \(a_2\rho\), con el orden de los
extremos ajustado al signo de \(a_2\).

\fi
\ifnum\HAFOderivationpart=9\relax
\subsection{Derivación completa de la identidad modal y su normalización}
\label{app:identidad_modal}

Sea \(P_0=P+bkr^{\T}\). Entonces
\[
\zeta I-P_0=(\zeta I-P)-bkr^{\T}.
\]
Si \(\zeta I-P\) es invertible, el lema del determinante matricial
\(\det(M+uv^{\T})=\det(M)(1+v^{\T}M^{-1}u)\), con
\[
M=\zeta I-P,\qquad u=-bk,\qquad v=r,
\]
produce
\begin{align}
\det(\zeta I-P_0)
&=\det(\zeta I-P)
\left[1-kr^{\T}(\zeta I-P)^{-1}b\right]\nonumber\\
&=\det(\zeta I-P)\bigl[1-kG(\zeta)\bigr],
\label{eq:lema_determinante_aplicado}
\end{align}
donde
\[
G(\zeta)=r^{\T}(\zeta I-P)^{-1}b,
\qquad W_q(s)=G(s^q).
\]
Por continuidad, la identidad polinómica se extiende también a los valores en
los que la resolvente de \(P\) es singular. Si
\[
1-kG(\zeta_0)=0,
\]
entonces \(\det(\zeta_0I-P_0)=0\), de modo que existe
\(v\neq0\) con \(P_0v=\zeta_0v\). Siempre que \(r^{\T}v\neq0\), se reemplaza
\(v\) por \(v/(r^{\T}v)\) y se obtiene \(r^{\T}v=1\). Ésta es la
normalización usada en \cref{eq:modo_normalizado,eq:semilla_modal}.

En la rama sesgada,
\[
X_1=\bigl[(\ii\omega_0)^qI-P\bigr]^{-1}b\Psi_1.
\]
Al premultiplicar por \(r^{\T}\),
\[
r^{\T}X_1
=W_q(\ii\omega_0)\Psi_1
=W_q(\ii\omega_0)N_1A.
\]
El cierre \(1-W_qN_1=0\) prueba algebraicamente que
\(r^{\T}X_1=A\).

\fi
\ifnum\HAFOderivationpart=10\relax
\subsection{Derivación completa de la homotopía afín}
\label{app:homotopia_afin}

La aproximación sesgada separa
\[
\psi(\sigma)
=\Psi_0+N_1(\sigma-c)
+\bigl[\psi(\sigma)-\Psi_0-N_1(\sigma-c)\bigr].
\]
Se definen
\begin{equation}
P_{\mathrm a}=P+N_1br^{\T},\qquad
d_{\mathrm a}=b(\Psi_0-N_1c).
\label{eq:parametros_homotopia_afin}
\end{equation}
La familia de campos es
\begin{equation}
\begin{aligned}
F_\varepsilon(X)
&=P_{\mathrm a}X+d_{\mathrm a}\\
&\quad+\varepsilon b
\left[
\psi(r^{\T}X)-\Psi_0-N_1(r^{\T}X-c)
\right],
\qquad 0\leq\varepsilon\leq1.
\end{aligned}
\label{eq:homotopia_afin}
\end{equation}

En \(\varepsilon=0\), la media \(\bar X\) de
\cref{eq:cierre_dc_sesgado} es un equilibrio del auxiliar:
\begin{align*}
F_0(\bar X)
&=P\bar X+bN_1c+b\Psi_0-bN_1c\\
&=P\bar X+b\Psi_0=0.
\end{align*}
Para una perturbación \(\xi=X-\bar X\),
\[
F_0(\bar X+\xi)=P_{\mathrm a}\xi,
\]
de modo que el auxiliar conserva tanto la media como el modo armónico.

En \(\varepsilon=1\), al agrupar términos,
\begin{align*}
F_1(X)
&=PX+bN_1r^{\T}X+b\Psi_0-bN_1c\\
&\quad+b\psi(r^{\T}X)-b\Psi_0-bN_1r^{\T}X+bN_1c\\
&=PX+b\psi(r^{\T}X).
\end{align*}
Así, el extremo final de \cref{eq:homotopia_afin} es exactamente el sistema
físico de \cref{eq:lure_comun}. Una homotopía centrada de la forma
\(P_0X+\varepsilon b[\psi(\sigma)-k\sigma]\) no conserva, en general, la
media \(c\neq0\); por eso no se utiliza para la rama sesgada.
\fi

\endgroup
\begingroup
\def\HAFOderivationpart{4}
% Este archivo es una biblioteca de derivaciones. Cada bloque se inserta en el
% lugar del cuerpo donde se usa, mediante \HAFOderivationpart. No se imprime
% como apéndice independiente.
\ifnum\HAFOderivationpart=1\relax
\subsection{Derivación completa de la forma de Lur'e}
\label{app:verificacion_lure}

Sea \(h(x)=\ell x+\psi(x)\). La primera ecuación de los dos modelos es
\[
{}^{\mathrm C}D_{t_0}^{q}x
=\alpha(y-x-h(x)).
\]
Al distribuir \(\alpha\),
\begin{align}
\alpha(y-x-h(x))
&=\alpha y-\alpha x-\alpha\ell x-\alpha\psi(x)\nonumber\\
&=-\alpha(1+\ell)x+\alpha y-\alpha\psi(x).
\label{eq:expansion_primera_fila}
\end{align}
Por otra parte,
\[
PX=
\begin{pmatrix}
-\alpha(1+\ell)x+\alpha y\\
x-y+z\\
-\beta y-\gamma z
\end{pmatrix},
\qquad
b\psi(r^{\T}X)=
\begin{pmatrix}
-\alpha\psi(x)\\0\\0
\end{pmatrix}.
\]
La suma reproduce \cref{eq:expansion_primera_fila} y las otras dos
ecuaciones. Para la saturación se toma
\(\ell=m_1\) y
\(\psi=(m_0-m_1)\operatorname{sat}\); para la arctangente,
\(\ell=a_1\) y \(\psi=a_2\arctan(\rho\,\cdot)\).

\fi
\ifnum\HAFOderivationpart=2\relax
\subsection{Derivación completa del determinante y la transferencia}
\label{app:transferencia_explicita}

Con \(\lambda=s^q\),
\begin{equation}
\lambda I-P=
\begin{pmatrix}
\lambda+\alpha(1+\ell)&-\alpha&0\\
-1&\lambda+1&-1\\
0&\beta&\lambda+\gamma
\end{pmatrix}.
\label{eq:matriz_resolvente}
\end{equation}
El menor de la posición \((1,1)\) es
\begin{align}
Q(\lambda)
&=
\det
\begin{pmatrix}
\lambda+1&-1\\
\beta&\lambda+\gamma
\end{pmatrix}\nonumber\\
&=(\lambda+1)(\lambda+\gamma)+\beta.
\label{eq:menor_Q}
\end{align}
Al desarrollar \(\det(\lambda I-P)\) por la primera fila,
\begin{align}
\Delta_\ell(\lambda)
&=\bigl[\lambda+\alpha(1+\ell)\bigr]Q(\lambda)
-(-\alpha)
\det
\begin{pmatrix}
-1&-1\\
0&\lambda+\gamma
\end{pmatrix}\nonumber\\
&=\bigl[\lambda+\alpha(1+\ell)\bigr]Q(\lambda)
-\alpha(\lambda+\gamma).
\label{eq:determinante_transferencia}
\end{align}

Para calcular la primera componente de
\((\lambda I-P)^{-1}b\), la regla de Cramer sustituye la primera columna de
\cref{eq:matriz_resolvente} por
\(b=(-\alpha,0,0)^{\T}\). El determinante del numerador es
\[
\det
\begin{pmatrix}
-\alpha&-\alpha&0\\
0&\lambda+1&-1\\
0&\beta&\lambda+\gamma
\end{pmatrix}
=-\alpha Q(\lambda).
\]
Como \(r^{\T}\) selecciona esa primera componente,
\begin{equation}
W_q(s)
=r^{\T}(s^qI-P)^{-1}b
=-\alpha\frac{Q(s^q)}{\Delta_\ell(s^q)}.
\label{eq:transferencia_cramer}
\end{equation}

La identidad
\[
P-s^qI=-(s^qI-P)
\]
implica
\[
r^{\T}(P-s^qI)^{-1}b=-W_q(s).
\]
Por ello, cambiar el orden de los términos dentro de la resolvente cambia
simultáneamente el signo de la transferencia y el signo del cierre. No cambia
la solución física si ambas modificaciones se realizan de manera coherente.

\fi
\ifnum\HAFOderivationpart=3\relax
\subsection{Derivación completa de los equilibrios}
\label{app:equilibrios}

En un equilibrio \(e=(x_e,y_e,z_e)^{\T}\), las dos últimas ecuaciones de
ambos sistemas dan
\[
0=x_e-y_e+z_e,\qquad
0=-\beta y_e-\gamma z_e.
\]
De la primera,
\[
z_e=y_e-x_e.
\]
Al sustituirla en la segunda,
\begin{align*}
0
&=-\beta y_e-\gamma(y_e-x_e)\\
&=-(\beta+\gamma)y_e+\gamma x_e.
\end{align*}
Si \(\beta+\gamma\neq0\),
\begin{equation}
y_e=\frac{\gamma}{\beta+\gamma}x_e,\qquad
z_e=-\frac{\beta}{\beta+\gamma}x_e.
\label{eq:equilibrios_yz}
\end{equation}
La primera ecuación de equilibrio exige
\[
0=y_e-x_e-h(x_e).
\]
Mediante \cref{eq:equilibrios_yz},
\begin{equation}
h(x_e)+\frac{\beta}{\beta+\gamma}x_e=0.
\label{eq:equilibrio_escalar_comun}
\end{equation}
Esta única ecuación escalar determina todos los equilibrios; las otras dos
coordenadas se recuperan después con \cref{eq:equilibrios_yz}.

\subsubsection{Equilibrios del Chua no suave}

En la región \(|x_e|<1\), \(h(x_e)=m_0x_e\). Entonces
\begin{equation}
\left(m_0+\frac{\beta}{\beta+\gamma}\right)x_e=0.
\label{eq:equilibrio_interior_ns}
\end{equation}
Cuando el coeficiente no es cero, la única raíz interior es \(x_e=0\).

Para \(x_e>1\),
\[
h(x_e)=m_1x_e+(m_0-m_1),
\]
y \cref{eq:equilibrio_escalar_comun} queda
\[
\left(m_1+\frac{\beta}{\beta+\gamma}\right)x_e+(m_0-m_1)=0.
\]
Por tanto,
\begin{equation}
x_\star=
\frac{m_1-m_0}
{m_1+\dfrac{\beta}{\beta+\gamma}}.
\label{eq:raiz_exterior_ns}
\end{equation}
Si \(x_\star>1\), existe el equilibrio positivo. Como la no linealidad es
impar, existe también la raíz \(-x_\star<-1\). Los tres equilibrios son
\begin{equation}
E_0=(0,0,0),\qquad
E_\pm=
\left(
\pm x_\star,\,
\pm\frac{\gamma}{\beta+\gamma}x_\star,\,
\mp\frac{\beta}{\beta+\gamma}x_\star
\right).
\label{eq:equilibrios_ns_simbolicos}
\end{equation}
La condición regional \(x_\star>1\) debe comprobarse; no basta con resolver la
ecuación lineal exterior.

\subsubsection{Equilibrios del Chua con arctangente}

Al usar
\(h(x)=a_1x+a_2\arctan(\rho x)\) en
\cref{eq:equilibrio_escalar_comun}, se obtiene
\begin{equation}
\left(a_1+\frac{\beta}{\beta+\gamma}\right)x_e
+a_2\arctan(\rho x_e)=0.
\label{eq:equilibrio_escalar_arctan}
\end{equation}
La expresión es impar, de modo que \(x_e=0\) siempre es raíz y toda raíz
positiva no nula tiene una compañera negativa. Las raíces no nulas se
determinan numéricamente y se sustituyen en \cref{eq:equilibrios_yz}. Esta
ecuación se evalúa con \(\rho=1\) para el caso bibliográfico y con el valor de
\(\rho\) de \cref{eq:parametros_c590} para el caso \texttt{c590}; no deben
intercambiarse las dos parametrizaciones.

\fi
\ifnum\HAFOderivationpart=4\relax
\subsection{Derivación completa de los jacobianos regionales}
\label{app:jacobianos}

Para la forma de Lur'e,
\[
F(X)=PX+b\psi(r^{\T}X).
\]
La regla de la cadena produce
\begin{equation}
J_e=P+b\,\psi'(r^{\T}e)r^{\T}.
\label{eq:jacobiano_lure}
\end{equation}
Como \(r^{\T}e=x_e\), solamente cambia la entrada \((1,1)\).

En el sistema no suave,
\[
\psi'_{\mathrm{ns}}(x)=
\begin{cases}
m_0-m_1,&|x|<1,\\
0,&|x|>1.
\end{cases}
\]
Por ello,
\begin{equation}
J_{\mathrm{in}}=
\begin{pmatrix}
-\alpha(1+m_0)&\alpha&0\\
1&-1&1\\
0&-\beta&-\gamma
\end{pmatrix},
\qquad
J_{\mathrm{out}}=
\begin{pmatrix}
-\alpha(1+m_1)&\alpha&0\\
1&-1&1\\
0&-\beta&-\gamma
\end{pmatrix}.
\label{eq:jacobianos_ns}
\end{equation}
En \(x=\pm1\) la derivada clásica no existe. Los equilibrios usados en este
reporte no se encuentran sobre esas superficies, por lo que se emplea el
jacobiano de la región que contiene a cada uno.

Para la arctangente,
\[
\psi'_{\mathrm a}(x)
=\frac{a_2\rho}{1+\rho^2x^2},
\]
y
\begin{equation}
J_e=
\begin{pmatrix}
\displaystyle
-\alpha\left(1+a_1+
\frac{a_2\rho}{1+\rho^2x_e^2}\right)
&\alpha&0\\[3mm]
1&-1&1\\
0&-\beta&-\gamma
\end{pmatrix}.
\label{eq:jacobiano_arctan}
\end{equation}
Los valores propios de
\cref{eq:jacobianos_ns,eq:jacobiano_arctan} se clasifican con
\cref{eq:matignon}; la estabilidad local no sustituye el sondeo de cuenca.

\fi
\ifnum\HAFOderivationpart=5\relax
\subsection{Derivación completa del polinomio característico}
\label{app:polinomio_caracteristico}

Sea \(g=h'(x_e)\) la pendiente total de la no linealidad en un equilibrio.
El jacobiano de cualquiera de las dos familias puede escribirse como
\begin{equation}
J(g)=
\begin{pmatrix}
-\alpha(1+g)&\alpha&0\\
1&-1&1\\
0&-\beta&-\gamma
\end{pmatrix}.
\label{eq:jacobiano_pendiente_total}
\end{equation}
Por tanto,
\[
\lambda I-J(g)=
\begin{pmatrix}
\lambda+\alpha(1+g)&-\alpha&0\\
-1&\lambda+1&-1\\
0&\beta&\lambda+\gamma
\end{pmatrix}.
\]
El desarrollo por la primera fila, igual al de
\cref{eq:determinante_transferencia}, produce
\begin{equation}
\chi_g(\lambda)
=\bigl[\lambda+\alpha(1+g)\bigr]
\bigl[(\lambda+1)(\lambda+\gamma)+\beta\bigr]
-\alpha(\lambda+\gamma).
\label{eq:polinomio_caracteristico_agrupado}
\end{equation}
Como
\[
(\lambda+1)(\lambda+\gamma)+\beta
=\lambda^2+(1+\gamma)\lambda+(\beta+\gamma),
\]
la multiplicación y agrupación por potencias de \(\lambda\) da
\begin{equation}
\begin{aligned}
\chi_g(\lambda)
={}&\lambda^3+
\bigl[1+\gamma+\alpha(1+g)\bigr]\lambda^2\\
&+\bigl[
\beta+\gamma+\alpha(1+g)(1+\gamma)-\alpha
\bigr]\lambda\\
&+\alpha\bigl[(1+g)(\beta+\gamma)-\gamma\bigr].
\end{aligned}
\label{eq:polinomio_caracteristico_expandido}
\end{equation}
En el sistema no suave se sustituye \(g=m_0\) en \(E_0\) y \(g=m_1\)
en \(E_\pm\). En el sistema con arctangente,
\[
g=a_1+\frac{a_2\rho}{1+\rho^2x_e^2}.
\]
Las raíces de \cref{eq:polinomio_caracteristico_expandido} son los valores
propios utilizados en \cref{eq:matignon}; así, la clasificación local se
obtiene a partir del mismo campo que define la integración.

\fi
\ifnum\HAFOderivationpart=6\relax
\subsection{Derivación completa para la saturación centrada}
\label{app:df_saturacion_centrada}

Sea
\[
\psi_{\mathrm{ns}}(\sigma)=\Delta\operatorname{sat}(\sigma),
\qquad \Delta=m_0-m_1.
\]
Si \(0<A\leq1\), no hay saturación y
\[
\psi_{\mathrm{ns}}(A\sin\theta)=\Delta A\sin\theta.
\]
Como \(\int_0^\pi\sin^2\theta\dd\theta=\pi/2\),
\[
N_{\mathrm{ns}}(A)
=\frac{2}{\pi A}\Delta A\frac{\pi}{2}
=\Delta.
\]

Para \(A>1\), defínase
\[
\theta_0=\arcsin\left(\frac1A\right),
\qquad 0<\theta_0<\frac{\pi}{2}.
\]
En \(0\leq\theta\leq\theta_0\) y
\(\pi-\theta_0\leq\theta\leq\pi\), la entrada permanece en la parte lineal;
en el intervalo central la saturación vale \(1\). Entonces
\begin{align}
N_{\mathrm{ns}}(A)
&=\frac{2\Delta}{\pi A}
\left[
2A\int_0^{\theta_0}\sin^2\theta\dd\theta
+\int_{\theta_0}^{\pi-\theta_0}\sin\theta\dd\theta
\right]\nonumber\\
&=\frac{2\Delta}{\pi A}
\left[
A\theta_0-\frac{A}{2}\sin(2\theta_0)
+2\cos\theta_0
\right]\nonumber\\
&=\frac{2\Delta}{\pi A}
\left[A\theta_0+\cos\theta_0\right].
\label{eq:df_sat_pasos}
\end{align}
En el último paso se usó
\[
\frac{A}{2}\sin(2\theta_0)
=A\sin\theta_0\cos\theta_0
=\cos\theta_0.
\]
Finalmente,
\[
\cos\theta_0
=\sqrt{1-\frac1{A^2}}
=\frac{\sqrt{A^2-1}}{A},
\]
y \cref{eq:df_sat_pasos} se convierte en
\[
N_{\mathrm{ns}}(A)
=\Delta\frac{2}{\pi}
\left[
\arcsin\left(\frac1A\right)
+\frac{\sqrt{A^2-1}}{A^2}
\right],
\]
que coincide con \cref{eq:df_saturacion}.

\fi
\ifnum\HAFOderivationpart=7\relax
\subsection{Derivación completa de la saturación sesgada}
\label{app:df_saturacion_sesgada}

Considérese
\[
u(\theta)=c+A\cos\theta,\qquad A>0.
\]
Para escribir una fórmula que incluya los casos parcialmente y completamente
saturados, se define
\[
[v]_{-1}^{1}=\max(-1,\min(1,v))
\]
y los ángulos
\begin{equation}
\theta_+
=\arccos\left[\frac{1-c}{A}\right]_{-1}^{1},
\qquad
\theta_-
=\arccos\left[\frac{-1-c}{A}\right]_{-1}^{1}.
\label{eq:angulos_saturacion_sesgada}
\end{equation}
Como \((1-c)/A>(-1-c)/A\) y \(\arccos\) es decreciente,
\(0\leq\theta_+\leq\theta_-\leq\pi\). En el intervalo
\([0,\pi]\),
\[
\operatorname{sat}(u(\theta))=
\begin{cases}
1,&0\leq\theta\leq\theta_+,\\
c+A\cos\theta,&\theta_+\leq\theta\leq\theta_-,\\
-1,&\theta_-\leq\theta\leq\pi.
\end{cases}
\]
La simetría respecto de \(2\pi\) permite integrar solamente en
\([0,\pi]\). Para \(\psi=\Delta\operatorname{sat}\),
\begin{align}
\Psi_0(c,A)
=\frac{\Delta}{\pi}\Big[
&\theta_+
+c(\theta_--\theta_+)
+A(\sin\theta_--\sin\theta_+)\nonumber\\
&-(\pi-\theta_-)
\Big].
\label{eq:psi0_saturacion_cerrada}
\end{align}

Para la componente fundamental se usan
\[
\int\cos\theta\dd\theta=\sin\theta,\qquad
\int\cos^2\theta\dd\theta
=\frac{\theta}{2}+\frac{\sin2\theta}{4}.
\]
La sustitución por intervalos da
\begin{align}
\Psi_1(c,A)
=\frac{2\Delta}{\pi}\Bigg[
&\sin\theta_++\sin\theta_-
+c(\sin\theta_--\sin\theta_+)\nonumber\\
&+\frac{A}{2}(\theta_--\theta_+)
+\frac{A}{4}
\bigl(\sin2\theta_--\sin2\theta_+\bigr)
\Bigg].
\label{eq:psi1_saturacion_cerrada}
\end{align}
Las operaciones de recorte en
\cref{eq:angulos_saturacion_sesgada} hacen que
\cref{eq:psi0_saturacion_cerrada,eq:psi1_saturacion_cerrada} sigan siendo
válidas cuando alguna de las tres regiones tiene longitud cero.

La ecuación de media puede reducirse a una igualdad escalar. De
\cref{eq:transferencia_cramer},
\begin{align}
W_q(0)
&=-\alpha
\frac{\beta+\gamma}
{\alpha(1+\ell)(\beta+\gamma)-\alpha\gamma}\nonumber\\
&=-\frac{\beta+\gamma}
{(1+\ell)(\beta+\gamma)-\gamma}.
\label{eq:ganancia_dc}
\end{align}
Por tanto,
\begin{equation}
c=W_q(0)\Psi_0(c,A),
\label{eq:cierre_dc_escalar}
\end{equation}
y el cierre armónico es
\[
1-W_q(\ii\omega)\frac{\Psi_1(c,A)}{A}=0.
\]
Estas tres ecuaciones reales ---una para la media, y las partes real e
imaginaria del cierre--- determinan \((c,A,\omega)\).

\fi
\ifnum\HAFOderivationpart=8\relax
\subsection{Derivación completa para la arctangente}
\label{app:df_arctan}

Sea
\[
\psi_{\mathrm a}(\sigma)=a_2\arctan(\rho\sigma),
\qquad
\xi=\rho A.
\]
Se define
\[
I(\xi)=
\int_0^\pi\arctan(\xi\sin\theta)\sin\theta\dd\theta.
\]
Como \(I(0)=0\), puede calcularse \(I\) integrando su derivada:
\begin{align}
I'(\xi)
&=\int_0^\pi
\frac{\sin^2\theta}{1+\xi^2\sin^2\theta}\dd\theta\nonumber\\
&=\frac1{\xi^2}
\left[
\pi-\int_0^\pi
\frac{\dd\theta}{1+\xi^2\sin^2\theta}
\right]\nonumber\\
&=\frac{\pi}{\xi^2}
\left(1-\frac1{\sqrt{1+\xi^2}}\right).
\label{eq:derivada_integral_arctan}
\end{align}
En la última igualdad se usó la integral estándar
\[
\int_0^\pi\frac{\dd\theta}{1+\xi^2\sin^2\theta}
=\frac{\pi}{\sqrt{1+\xi^2}}.
\]
Una primitiva que se anula en \(\xi=0\) es
\begin{equation}
I(\xi)
=\pi\frac{\sqrt{1+\xi^2}-1}{\xi}.
\label{eq:integral_arctan}
\end{equation}
Al sustituir \(\xi=\rho A\) en
\cref{eq:df_centrada},
\begin{align}
N_{\mathrm a}(A)
&=\frac{2a_2}{\pi A}I(\rho A)\nonumber\\
&=\frac{2a_2}{\rho A^2}
\left(\sqrt{1+\rho^2A^2}-1\right)\nonumber\\
&=\frac{2a_2\rho}
{1+\sqrt{1+\rho^2A^2}}.
\label{eq:df_arctan_derivada}
\end{align}

Si el cierre exige \(N_{\mathrm a}(A)=k\), la amplitud puede despejarse.
De la forma racionalizada,
\[
1+\sqrt{1+\rho^2A^2}=\frac{2a_2\rho}{k}.
\]
Al aislar la raíz, elevar al cuadrado y simplificar,
\begin{equation}
A^2
=\frac{4a_2(a_2\rho-k)}{k^2\rho}.
\label{eq:amplitud_arctan}
\end{equation}
La solución se acepta solamente si el lado derecho es positivo y la igualdad
previa a elevar al cuadrado tiene lado derecho no menor que \(2\). De manera
equivalente, para \(A>0\), \(k\) tiene el mismo signo que \(a_2\) y se
encuentra estrictamente entre \(0\) y \(a_2\rho\), con el orden de los
extremos ajustado al signo de \(a_2\).

\fi
\ifnum\HAFOderivationpart=9\relax
\subsection{Derivación completa de la identidad modal y su normalización}
\label{app:identidad_modal}

Sea \(P_0=P+bkr^{\T}\). Entonces
\[
\zeta I-P_0=(\zeta I-P)-bkr^{\T}.
\]
Si \(\zeta I-P\) es invertible, el lema del determinante matricial
\(\det(M+uv^{\T})=\det(M)(1+v^{\T}M^{-1}u)\), con
\[
M=\zeta I-P,\qquad u=-bk,\qquad v=r,
\]
produce
\begin{align}
\det(\zeta I-P_0)
&=\det(\zeta I-P)
\left[1-kr^{\T}(\zeta I-P)^{-1}b\right]\nonumber\\
&=\det(\zeta I-P)\bigl[1-kG(\zeta)\bigr],
\label{eq:lema_determinante_aplicado}
\end{align}
donde
\[
G(\zeta)=r^{\T}(\zeta I-P)^{-1}b,
\qquad W_q(s)=G(s^q).
\]
Por continuidad, la identidad polinómica se extiende también a los valores en
los que la resolvente de \(P\) es singular. Si
\[
1-kG(\zeta_0)=0,
\]
entonces \(\det(\zeta_0I-P_0)=0\), de modo que existe
\(v\neq0\) con \(P_0v=\zeta_0v\). Siempre que \(r^{\T}v\neq0\), se reemplaza
\(v\) por \(v/(r^{\T}v)\) y se obtiene \(r^{\T}v=1\). Ésta es la
normalización usada en \cref{eq:modo_normalizado,eq:semilla_modal}.

En la rama sesgada,
\[
X_1=\bigl[(\ii\omega_0)^qI-P\bigr]^{-1}b\Psi_1.
\]
Al premultiplicar por \(r^{\T}\),
\[
r^{\T}X_1
=W_q(\ii\omega_0)\Psi_1
=W_q(\ii\omega_0)N_1A.
\]
El cierre \(1-W_qN_1=0\) prueba algebraicamente que
\(r^{\T}X_1=A\).

\fi
\ifnum\HAFOderivationpart=10\relax
\subsection{Derivación completa de la homotopía afín}
\label{app:homotopia_afin}

La aproximación sesgada separa
\[
\psi(\sigma)
=\Psi_0+N_1(\sigma-c)
+\bigl[\psi(\sigma)-\Psi_0-N_1(\sigma-c)\bigr].
\]
Se definen
\begin{equation}
P_{\mathrm a}=P+N_1br^{\T},\qquad
d_{\mathrm a}=b(\Psi_0-N_1c).
\label{eq:parametros_homotopia_afin}
\end{equation}
La familia de campos es
\begin{equation}
\begin{aligned}
F_\varepsilon(X)
&=P_{\mathrm a}X+d_{\mathrm a}\\
&\quad+\varepsilon b
\left[
\psi(r^{\T}X)-\Psi_0-N_1(r^{\T}X-c)
\right],
\qquad 0\leq\varepsilon\leq1.
\end{aligned}
\label{eq:homotopia_afin}
\end{equation}

En \(\varepsilon=0\), la media \(\bar X\) de
\cref{eq:cierre_dc_sesgado} es un equilibrio del auxiliar:
\begin{align*}
F_0(\bar X)
&=P\bar X+bN_1c+b\Psi_0-bN_1c\\
&=P\bar X+b\Psi_0=0.
\end{align*}
Para una perturbación \(\xi=X-\bar X\),
\[
F_0(\bar X+\xi)=P_{\mathrm a}\xi,
\]
de modo que el auxiliar conserva tanto la media como el modo armónico.

En \(\varepsilon=1\), al agrupar términos,
\begin{align*}
F_1(X)
&=PX+bN_1r^{\T}X+b\Psi_0-bN_1c\\
&\quad+b\psi(r^{\T}X)-b\Psi_0-bN_1r^{\T}X+bN_1c\\
&=PX+b\psi(r^{\T}X).
\end{align*}
Así, el extremo final de \cref{eq:homotopia_afin} es exactamente el sistema
físico de \cref{eq:lure_comun}. Una homotopía centrada de la forma
\(P_0X+\varepsilon b[\psi(\sigma)-k\sigma]\) no conserva, en general, la
media \(c\neq0\); por eso no se utiliza para la rama sesgada.
\fi

\endgroup
\begingroup
\def\HAFOderivationpart{5}
% Este archivo es una biblioteca de derivaciones. Cada bloque se inserta en el
% lugar del cuerpo donde se usa, mediante \HAFOderivationpart. No se imprime
% como apéndice independiente.
\ifnum\HAFOderivationpart=1\relax
\subsection{Derivación completa de la forma de Lur'e}
\label{app:verificacion_lure}

Sea \(h(x)=\ell x+\psi(x)\). La primera ecuación de los dos modelos es
\[
{}^{\mathrm C}D_{t_0}^{q}x
=\alpha(y-x-h(x)).
\]
Al distribuir \(\alpha\),
\begin{align}
\alpha(y-x-h(x))
&=\alpha y-\alpha x-\alpha\ell x-\alpha\psi(x)\nonumber\\
&=-\alpha(1+\ell)x+\alpha y-\alpha\psi(x).
\label{eq:expansion_primera_fila}
\end{align}
Por otra parte,
\[
PX=
\begin{pmatrix}
-\alpha(1+\ell)x+\alpha y\\
x-y+z\\
-\beta y-\gamma z
\end{pmatrix},
\qquad
b\psi(r^{\T}X)=
\begin{pmatrix}
-\alpha\psi(x)\\0\\0
\end{pmatrix}.
\]
La suma reproduce \cref{eq:expansion_primera_fila} y las otras dos
ecuaciones. Para la saturación se toma
\(\ell=m_1\) y
\(\psi=(m_0-m_1)\operatorname{sat}\); para la arctangente,
\(\ell=a_1\) y \(\psi=a_2\arctan(\rho\,\cdot)\).

\fi
\ifnum\HAFOderivationpart=2\relax
\subsection{Derivación completa del determinante y la transferencia}
\label{app:transferencia_explicita}

Con \(\lambda=s^q\),
\begin{equation}
\lambda I-P=
\begin{pmatrix}
\lambda+\alpha(1+\ell)&-\alpha&0\\
-1&\lambda+1&-1\\
0&\beta&\lambda+\gamma
\end{pmatrix}.
\label{eq:matriz_resolvente}
\end{equation}
El menor de la posición \((1,1)\) es
\begin{align}
Q(\lambda)
&=
\det
\begin{pmatrix}
\lambda+1&-1\\
\beta&\lambda+\gamma
\end{pmatrix}\nonumber\\
&=(\lambda+1)(\lambda+\gamma)+\beta.
\label{eq:menor_Q}
\end{align}
Al desarrollar \(\det(\lambda I-P)\) por la primera fila,
\begin{align}
\Delta_\ell(\lambda)
&=\bigl[\lambda+\alpha(1+\ell)\bigr]Q(\lambda)
-(-\alpha)
\det
\begin{pmatrix}
-1&-1\\
0&\lambda+\gamma
\end{pmatrix}\nonumber\\
&=\bigl[\lambda+\alpha(1+\ell)\bigr]Q(\lambda)
-\alpha(\lambda+\gamma).
\label{eq:determinante_transferencia}
\end{align}

Para calcular la primera componente de
\((\lambda I-P)^{-1}b\), la regla de Cramer sustituye la primera columna de
\cref{eq:matriz_resolvente} por
\(b=(-\alpha,0,0)^{\T}\). El determinante del numerador es
\[
\det
\begin{pmatrix}
-\alpha&-\alpha&0\\
0&\lambda+1&-1\\
0&\beta&\lambda+\gamma
\end{pmatrix}
=-\alpha Q(\lambda).
\]
Como \(r^{\T}\) selecciona esa primera componente,
\begin{equation}
W_q(s)
=r^{\T}(s^qI-P)^{-1}b
=-\alpha\frac{Q(s^q)}{\Delta_\ell(s^q)}.
\label{eq:transferencia_cramer}
\end{equation}

La identidad
\[
P-s^qI=-(s^qI-P)
\]
implica
\[
r^{\T}(P-s^qI)^{-1}b=-W_q(s).
\]
Por ello, cambiar el orden de los términos dentro de la resolvente cambia
simultáneamente el signo de la transferencia y el signo del cierre. No cambia
la solución física si ambas modificaciones se realizan de manera coherente.

\fi
\ifnum\HAFOderivationpart=3\relax
\subsection{Derivación completa de los equilibrios}
\label{app:equilibrios}

En un equilibrio \(e=(x_e,y_e,z_e)^{\T}\), las dos últimas ecuaciones de
ambos sistemas dan
\[
0=x_e-y_e+z_e,\qquad
0=-\beta y_e-\gamma z_e.
\]
De la primera,
\[
z_e=y_e-x_e.
\]
Al sustituirla en la segunda,
\begin{align*}
0
&=-\beta y_e-\gamma(y_e-x_e)\\
&=-(\beta+\gamma)y_e+\gamma x_e.
\end{align*}
Si \(\beta+\gamma\neq0\),
\begin{equation}
y_e=\frac{\gamma}{\beta+\gamma}x_e,\qquad
z_e=-\frac{\beta}{\beta+\gamma}x_e.
\label{eq:equilibrios_yz}
\end{equation}
La primera ecuación de equilibrio exige
\[
0=y_e-x_e-h(x_e).
\]
Mediante \cref{eq:equilibrios_yz},
\begin{equation}
h(x_e)+\frac{\beta}{\beta+\gamma}x_e=0.
\label{eq:equilibrio_escalar_comun}
\end{equation}
Esta única ecuación escalar determina todos los equilibrios; las otras dos
coordenadas se recuperan después con \cref{eq:equilibrios_yz}.

\subsubsection{Equilibrios del Chua no suave}

En la región \(|x_e|<1\), \(h(x_e)=m_0x_e\). Entonces
\begin{equation}
\left(m_0+\frac{\beta}{\beta+\gamma}\right)x_e=0.
\label{eq:equilibrio_interior_ns}
\end{equation}
Cuando el coeficiente no es cero, la única raíz interior es \(x_e=0\).

Para \(x_e>1\),
\[
h(x_e)=m_1x_e+(m_0-m_1),
\]
y \cref{eq:equilibrio_escalar_comun} queda
\[
\left(m_1+\frac{\beta}{\beta+\gamma}\right)x_e+(m_0-m_1)=0.
\]
Por tanto,
\begin{equation}
x_\star=
\frac{m_1-m_0}
{m_1+\dfrac{\beta}{\beta+\gamma}}.
\label{eq:raiz_exterior_ns}
\end{equation}
Si \(x_\star>1\), existe el equilibrio positivo. Como la no linealidad es
impar, existe también la raíz \(-x_\star<-1\). Los tres equilibrios son
\begin{equation}
E_0=(0,0,0),\qquad
E_\pm=
\left(
\pm x_\star,\,
\pm\frac{\gamma}{\beta+\gamma}x_\star,\,
\mp\frac{\beta}{\beta+\gamma}x_\star
\right).
\label{eq:equilibrios_ns_simbolicos}
\end{equation}
La condición regional \(x_\star>1\) debe comprobarse; no basta con resolver la
ecuación lineal exterior.

\subsubsection{Equilibrios del Chua con arctangente}

Al usar
\(h(x)=a_1x+a_2\arctan(\rho x)\) en
\cref{eq:equilibrio_escalar_comun}, se obtiene
\begin{equation}
\left(a_1+\frac{\beta}{\beta+\gamma}\right)x_e
+a_2\arctan(\rho x_e)=0.
\label{eq:equilibrio_escalar_arctan}
\end{equation}
La expresión es impar, de modo que \(x_e=0\) siempre es raíz y toda raíz
positiva no nula tiene una compañera negativa. Las raíces no nulas se
determinan numéricamente y se sustituyen en \cref{eq:equilibrios_yz}. Esta
ecuación se evalúa con \(\rho=1\) para el caso bibliográfico y con el valor de
\(\rho\) de \cref{eq:parametros_c590} para el caso \texttt{c590}; no deben
intercambiarse las dos parametrizaciones.

\fi
\ifnum\HAFOderivationpart=4\relax
\subsection{Derivación completa de los jacobianos regionales}
\label{app:jacobianos}

Para la forma de Lur'e,
\[
F(X)=PX+b\psi(r^{\T}X).
\]
La regla de la cadena produce
\begin{equation}
J_e=P+b\,\psi'(r^{\T}e)r^{\T}.
\label{eq:jacobiano_lure}
\end{equation}
Como \(r^{\T}e=x_e\), solamente cambia la entrada \((1,1)\).

En el sistema no suave,
\[
\psi'_{\mathrm{ns}}(x)=
\begin{cases}
m_0-m_1,&|x|<1,\\
0,&|x|>1.
\end{cases}
\]
Por ello,
\begin{equation}
J_{\mathrm{in}}=
\begin{pmatrix}
-\alpha(1+m_0)&\alpha&0\\
1&-1&1\\
0&-\beta&-\gamma
\end{pmatrix},
\qquad
J_{\mathrm{out}}=
\begin{pmatrix}
-\alpha(1+m_1)&\alpha&0\\
1&-1&1\\
0&-\beta&-\gamma
\end{pmatrix}.
\label{eq:jacobianos_ns}
\end{equation}
En \(x=\pm1\) la derivada clásica no existe. Los equilibrios usados en este
reporte no se encuentran sobre esas superficies, por lo que se emplea el
jacobiano de la región que contiene a cada uno.

Para la arctangente,
\[
\psi'_{\mathrm a}(x)
=\frac{a_2\rho}{1+\rho^2x^2},
\]
y
\begin{equation}
J_e=
\begin{pmatrix}
\displaystyle
-\alpha\left(1+a_1+
\frac{a_2\rho}{1+\rho^2x_e^2}\right)
&\alpha&0\\[3mm]
1&-1&1\\
0&-\beta&-\gamma
\end{pmatrix}.
\label{eq:jacobiano_arctan}
\end{equation}
Los valores propios de
\cref{eq:jacobianos_ns,eq:jacobiano_arctan} se clasifican con
\cref{eq:matignon}; la estabilidad local no sustituye el sondeo de cuenca.

\fi
\ifnum\HAFOderivationpart=5\relax
\subsection{Derivación completa del polinomio característico}
\label{app:polinomio_caracteristico}

Sea \(g=h'(x_e)\) la pendiente total de la no linealidad en un equilibrio.
El jacobiano de cualquiera de las dos familias puede escribirse como
\begin{equation}
J(g)=
\begin{pmatrix}
-\alpha(1+g)&\alpha&0\\
1&-1&1\\
0&-\beta&-\gamma
\end{pmatrix}.
\label{eq:jacobiano_pendiente_total}
\end{equation}
Por tanto,
\[
\lambda I-J(g)=
\begin{pmatrix}
\lambda+\alpha(1+g)&-\alpha&0\\
-1&\lambda+1&-1\\
0&\beta&\lambda+\gamma
\end{pmatrix}.
\]
El desarrollo por la primera fila, igual al de
\cref{eq:determinante_transferencia}, produce
\begin{equation}
\chi_g(\lambda)
=\bigl[\lambda+\alpha(1+g)\bigr]
\bigl[(\lambda+1)(\lambda+\gamma)+\beta\bigr]
-\alpha(\lambda+\gamma).
\label{eq:polinomio_caracteristico_agrupado}
\end{equation}
Como
\[
(\lambda+1)(\lambda+\gamma)+\beta
=\lambda^2+(1+\gamma)\lambda+(\beta+\gamma),
\]
la multiplicación y agrupación por potencias de \(\lambda\) da
\begin{equation}
\begin{aligned}
\chi_g(\lambda)
={}&\lambda^3+
\bigl[1+\gamma+\alpha(1+g)\bigr]\lambda^2\\
&+\bigl[
\beta+\gamma+\alpha(1+g)(1+\gamma)-\alpha
\bigr]\lambda\\
&+\alpha\bigl[(1+g)(\beta+\gamma)-\gamma\bigr].
\end{aligned}
\label{eq:polinomio_caracteristico_expandido}
\end{equation}
En el sistema no suave se sustituye \(g=m_0\) en \(E_0\) y \(g=m_1\)
en \(E_\pm\). En el sistema con arctangente,
\[
g=a_1+\frac{a_2\rho}{1+\rho^2x_e^2}.
\]
Las raíces de \cref{eq:polinomio_caracteristico_expandido} son los valores
propios utilizados en \cref{eq:matignon}; así, la clasificación local se
obtiene a partir del mismo campo que define la integración.

\fi
\ifnum\HAFOderivationpart=6\relax
\subsection{Derivación completa para la saturación centrada}
\label{app:df_saturacion_centrada}

Sea
\[
\psi_{\mathrm{ns}}(\sigma)=\Delta\operatorname{sat}(\sigma),
\qquad \Delta=m_0-m_1.
\]
Si \(0<A\leq1\), no hay saturación y
\[
\psi_{\mathrm{ns}}(A\sin\theta)=\Delta A\sin\theta.
\]
Como \(\int_0^\pi\sin^2\theta\dd\theta=\pi/2\),
\[
N_{\mathrm{ns}}(A)
=\frac{2}{\pi A}\Delta A\frac{\pi}{2}
=\Delta.
\]

Para \(A>1\), defínase
\[
\theta_0=\arcsin\left(\frac1A\right),
\qquad 0<\theta_0<\frac{\pi}{2}.
\]
En \(0\leq\theta\leq\theta_0\) y
\(\pi-\theta_0\leq\theta\leq\pi\), la entrada permanece en la parte lineal;
en el intervalo central la saturación vale \(1\). Entonces
\begin{align}
N_{\mathrm{ns}}(A)
&=\frac{2\Delta}{\pi A}
\left[
2A\int_0^{\theta_0}\sin^2\theta\dd\theta
+\int_{\theta_0}^{\pi-\theta_0}\sin\theta\dd\theta
\right]\nonumber\\
&=\frac{2\Delta}{\pi A}
\left[
A\theta_0-\frac{A}{2}\sin(2\theta_0)
+2\cos\theta_0
\right]\nonumber\\
&=\frac{2\Delta}{\pi A}
\left[A\theta_0+\cos\theta_0\right].
\label{eq:df_sat_pasos}
\end{align}
En el último paso se usó
\[
\frac{A}{2}\sin(2\theta_0)
=A\sin\theta_0\cos\theta_0
=\cos\theta_0.
\]
Finalmente,
\[
\cos\theta_0
=\sqrt{1-\frac1{A^2}}
=\frac{\sqrt{A^2-1}}{A},
\]
y \cref{eq:df_sat_pasos} se convierte en
\[
N_{\mathrm{ns}}(A)
=\Delta\frac{2}{\pi}
\left[
\arcsin\left(\frac1A\right)
+\frac{\sqrt{A^2-1}}{A^2}
\right],
\]
que coincide con \cref{eq:df_saturacion}.

\fi
\ifnum\HAFOderivationpart=7\relax
\subsection{Derivación completa de la saturación sesgada}
\label{app:df_saturacion_sesgada}

Considérese
\[
u(\theta)=c+A\cos\theta,\qquad A>0.
\]
Para escribir una fórmula que incluya los casos parcialmente y completamente
saturados, se define
\[
[v]_{-1}^{1}=\max(-1,\min(1,v))
\]
y los ángulos
\begin{equation}
\theta_+
=\arccos\left[\frac{1-c}{A}\right]_{-1}^{1},
\qquad
\theta_-
=\arccos\left[\frac{-1-c}{A}\right]_{-1}^{1}.
\label{eq:angulos_saturacion_sesgada}
\end{equation}
Como \((1-c)/A>(-1-c)/A\) y \(\arccos\) es decreciente,
\(0\leq\theta_+\leq\theta_-\leq\pi\). En el intervalo
\([0,\pi]\),
\[
\operatorname{sat}(u(\theta))=
\begin{cases}
1,&0\leq\theta\leq\theta_+,\\
c+A\cos\theta,&\theta_+\leq\theta\leq\theta_-,\\
-1,&\theta_-\leq\theta\leq\pi.
\end{cases}
\]
La simetría respecto de \(2\pi\) permite integrar solamente en
\([0,\pi]\). Para \(\psi=\Delta\operatorname{sat}\),
\begin{align}
\Psi_0(c,A)
=\frac{\Delta}{\pi}\Big[
&\theta_+
+c(\theta_--\theta_+)
+A(\sin\theta_--\sin\theta_+)\nonumber\\
&-(\pi-\theta_-)
\Big].
\label{eq:psi0_saturacion_cerrada}
\end{align}

Para la componente fundamental se usan
\[
\int\cos\theta\dd\theta=\sin\theta,\qquad
\int\cos^2\theta\dd\theta
=\frac{\theta}{2}+\frac{\sin2\theta}{4}.
\]
La sustitución por intervalos da
\begin{align}
\Psi_1(c,A)
=\frac{2\Delta}{\pi}\Bigg[
&\sin\theta_++\sin\theta_-
+c(\sin\theta_--\sin\theta_+)\nonumber\\
&+\frac{A}{2}(\theta_--\theta_+)
+\frac{A}{4}
\bigl(\sin2\theta_--\sin2\theta_+\bigr)
\Bigg].
\label{eq:psi1_saturacion_cerrada}
\end{align}
Las operaciones de recorte en
\cref{eq:angulos_saturacion_sesgada} hacen que
\cref{eq:psi0_saturacion_cerrada,eq:psi1_saturacion_cerrada} sigan siendo
válidas cuando alguna de las tres regiones tiene longitud cero.

La ecuación de media puede reducirse a una igualdad escalar. De
\cref{eq:transferencia_cramer},
\begin{align}
W_q(0)
&=-\alpha
\frac{\beta+\gamma}
{\alpha(1+\ell)(\beta+\gamma)-\alpha\gamma}\nonumber\\
&=-\frac{\beta+\gamma}
{(1+\ell)(\beta+\gamma)-\gamma}.
\label{eq:ganancia_dc}
\end{align}
Por tanto,
\begin{equation}
c=W_q(0)\Psi_0(c,A),
\label{eq:cierre_dc_escalar}
\end{equation}
y el cierre armónico es
\[
1-W_q(\ii\omega)\frac{\Psi_1(c,A)}{A}=0.
\]
Estas tres ecuaciones reales ---una para la media, y las partes real e
imaginaria del cierre--- determinan \((c,A,\omega)\).

\fi
\ifnum\HAFOderivationpart=8\relax
\subsection{Derivación completa para la arctangente}
\label{app:df_arctan}

Sea
\[
\psi_{\mathrm a}(\sigma)=a_2\arctan(\rho\sigma),
\qquad
\xi=\rho A.
\]
Se define
\[
I(\xi)=
\int_0^\pi\arctan(\xi\sin\theta)\sin\theta\dd\theta.
\]
Como \(I(0)=0\), puede calcularse \(I\) integrando su derivada:
\begin{align}
I'(\xi)
&=\int_0^\pi
\frac{\sin^2\theta}{1+\xi^2\sin^2\theta}\dd\theta\nonumber\\
&=\frac1{\xi^2}
\left[
\pi-\int_0^\pi
\frac{\dd\theta}{1+\xi^2\sin^2\theta}
\right]\nonumber\\
&=\frac{\pi}{\xi^2}
\left(1-\frac1{\sqrt{1+\xi^2}}\right).
\label{eq:derivada_integral_arctan}
\end{align}
En la última igualdad se usó la integral estándar
\[
\int_0^\pi\frac{\dd\theta}{1+\xi^2\sin^2\theta}
=\frac{\pi}{\sqrt{1+\xi^2}}.
\]
Una primitiva que se anula en \(\xi=0\) es
\begin{equation}
I(\xi)
=\pi\frac{\sqrt{1+\xi^2}-1}{\xi}.
\label{eq:integral_arctan}
\end{equation}
Al sustituir \(\xi=\rho A\) en
\cref{eq:df_centrada},
\begin{align}
N_{\mathrm a}(A)
&=\frac{2a_2}{\pi A}I(\rho A)\nonumber\\
&=\frac{2a_2}{\rho A^2}
\left(\sqrt{1+\rho^2A^2}-1\right)\nonumber\\
&=\frac{2a_2\rho}
{1+\sqrt{1+\rho^2A^2}}.
\label{eq:df_arctan_derivada}
\end{align}

Si el cierre exige \(N_{\mathrm a}(A)=k\), la amplitud puede despejarse.
De la forma racionalizada,
\[
1+\sqrt{1+\rho^2A^2}=\frac{2a_2\rho}{k}.
\]
Al aislar la raíz, elevar al cuadrado y simplificar,
\begin{equation}
A^2
=\frac{4a_2(a_2\rho-k)}{k^2\rho}.
\label{eq:amplitud_arctan}
\end{equation}
La solución se acepta solamente si el lado derecho es positivo y la igualdad
previa a elevar al cuadrado tiene lado derecho no menor que \(2\). De manera
equivalente, para \(A>0\), \(k\) tiene el mismo signo que \(a_2\) y se
encuentra estrictamente entre \(0\) y \(a_2\rho\), con el orden de los
extremos ajustado al signo de \(a_2\).

\fi
\ifnum\HAFOderivationpart=9\relax
\subsection{Derivación completa de la identidad modal y su normalización}
\label{app:identidad_modal}

Sea \(P_0=P+bkr^{\T}\). Entonces
\[
\zeta I-P_0=(\zeta I-P)-bkr^{\T}.
\]
Si \(\zeta I-P\) es invertible, el lema del determinante matricial
\(\det(M+uv^{\T})=\det(M)(1+v^{\T}M^{-1}u)\), con
\[
M=\zeta I-P,\qquad u=-bk,\qquad v=r,
\]
produce
\begin{align}
\det(\zeta I-P_0)
&=\det(\zeta I-P)
\left[1-kr^{\T}(\zeta I-P)^{-1}b\right]\nonumber\\
&=\det(\zeta I-P)\bigl[1-kG(\zeta)\bigr],
\label{eq:lema_determinante_aplicado}
\end{align}
donde
\[
G(\zeta)=r^{\T}(\zeta I-P)^{-1}b,
\qquad W_q(s)=G(s^q).
\]
Por continuidad, la identidad polinómica se extiende también a los valores en
los que la resolvente de \(P\) es singular. Si
\[
1-kG(\zeta_0)=0,
\]
entonces \(\det(\zeta_0I-P_0)=0\), de modo que existe
\(v\neq0\) con \(P_0v=\zeta_0v\). Siempre que \(r^{\T}v\neq0\), se reemplaza
\(v\) por \(v/(r^{\T}v)\) y se obtiene \(r^{\T}v=1\). Ésta es la
normalización usada en \cref{eq:modo_normalizado,eq:semilla_modal}.

En la rama sesgada,
\[
X_1=\bigl[(\ii\omega_0)^qI-P\bigr]^{-1}b\Psi_1.
\]
Al premultiplicar por \(r^{\T}\),
\[
r^{\T}X_1
=W_q(\ii\omega_0)\Psi_1
=W_q(\ii\omega_0)N_1A.
\]
El cierre \(1-W_qN_1=0\) prueba algebraicamente que
\(r^{\T}X_1=A\).

\fi
\ifnum\HAFOderivationpart=10\relax
\subsection{Derivación completa de la homotopía afín}
\label{app:homotopia_afin}

La aproximación sesgada separa
\[
\psi(\sigma)
=\Psi_0+N_1(\sigma-c)
+\bigl[\psi(\sigma)-\Psi_0-N_1(\sigma-c)\bigr].
\]
Se definen
\begin{equation}
P_{\mathrm a}=P+N_1br^{\T},\qquad
d_{\mathrm a}=b(\Psi_0-N_1c).
\label{eq:parametros_homotopia_afin}
\end{equation}
La familia de campos es
\begin{equation}
\begin{aligned}
F_\varepsilon(X)
&=P_{\mathrm a}X+d_{\mathrm a}\\
&\quad+\varepsilon b
\left[
\psi(r^{\T}X)-\Psi_0-N_1(r^{\T}X-c)
\right],
\qquad 0\leq\varepsilon\leq1.
\end{aligned}
\label{eq:homotopia_afin}
\end{equation}

En \(\varepsilon=0\), la media \(\bar X\) de
\cref{eq:cierre_dc_sesgado} es un equilibrio del auxiliar:
\begin{align*}
F_0(\bar X)
&=P\bar X+bN_1c+b\Psi_0-bN_1c\\
&=P\bar X+b\Psi_0=0.
\end{align*}
Para una perturbación \(\xi=X-\bar X\),
\[
F_0(\bar X+\xi)=P_{\mathrm a}\xi,
\]
de modo que el auxiliar conserva tanto la media como el modo armónico.

En \(\varepsilon=1\), al agrupar términos,
\begin{align*}
F_1(X)
&=PX+bN_1r^{\T}X+b\Psi_0-bN_1c\\
&\quad+b\psi(r^{\T}X)-b\Psi_0-bN_1r^{\T}X+bN_1c\\
&=PX+b\psi(r^{\T}X).
\end{align*}
Así, el extremo final de \cref{eq:homotopia_afin} es exactamente el sistema
físico de \cref{eq:lure_comun}. Una homotopía centrada de la forma
\(P_0X+\varepsilon b[\psi(\sigma)-k\sigma]\) no conserva, en general, la
media \(c\neq0\); por eso no se utiliza para la rama sesgada.
\fi

\endgroup
\fi

\ifdefined\HAFOModelsOnly\else
\begingroup
% Sin selector se conserva el bloque completo. En la edición de estudio,
% HAFOFractionalBlockSelect permite insertar por separado los operadores y la
% estabilidad sin duplicar el contenido.
\ifdefined\HAFOFractionalBlockSelect\else
  \def\HAFOFractionalBlockOperators{1}
  \def\HAFOFractionalBlockStability{1}
\fi

\ifdefined\HAFOFractionalBlockOperators
\section{Fundamentos fraccionarios}
\label{sec:fundamentos}

\HAFOCardFractional

\subsection{Funciones especiales e integral de Riemann--Liouville}

El cálculo fraccionario se apoya en tres objetos previos a la definición de una
derivada. Para \(z>0\) y \(a,b>0\),
\begin{align}
 \Gamma(z)&=\int_0^\infty t^{z-1}e^{-t}\,\dd t,
 &\Gamma(z+1)&=z\Gamma(z),\label{eq:gamma-fractional}\\
 B(a,b)&=\int_0^1 u^{a-1}(1-u)^{b-1}\,\dd u
 =\frac{\Gamma(a)\Gamma(b)}{\Gamma(a+b)}.\label{eq:beta-gamma}
\end{align}
La función de Mittag--Leffler de dos parámetros,
\begin{equation}
 E_{\alpha,\beta}(z)=\sum_{k=0}^{\infty}
 \frac{z^k}{\Gamma(\alpha k+\beta)},
 \qquad \alpha>0,
 \label{eq:mittag-leffler-definition}
\end{equation}
generaliza a la exponencial porque \(E_{1,1}(z)=e^z\). Estas funciones
aparecen de manera natural al resolver ecuaciones lineales fraccionarias
\cite{Podlubny1999,Monje2010,Oliveira2019SolvedFractional}.

Para \(\alpha>0\), la integral fraccionaria izquierda de
Riemann--Liouville es
\begin{equation}
 (I_{t_0}^{\alpha}g)(t)=\frac{1}{\Gamma(\alpha)}
 \int_{t_0}^{t}(t-\tau)^{\alpha-1}g(\tau)\,\dd\tau.
 \label{eq:rl-integral}
\end{equation}
Para \(g\) continua en un intervalo compacto, Fubini permite invertir el
orden de integración en \(I_{t_0}^{\alpha}I_{t_0}^{\beta}g\). La integral
interior resultante se evalúa con \(\tau=s+(t-s)u\):
\begin{align*}
 \int_s^t(t-\tau)^{\alpha-1}(\tau-s)^{\beta-1}\,\mathrm d\tau
 &=(t-s)^{\alpha+\beta-1}
   \int_0^1(1-u)^{\alpha-1}u^{\beta-1}\,\mathrm du\\
 &=(t-s)^{\alpha+\beta-1}
   \frac{\Gamma(\alpha)\Gamma(\beta)}{\Gamma(\alpha+\beta)}.
\end{align*}
Al cancelar los factores Gamma exteriores se obtiene
\(I_{t_0}^{\alpha}I_{t_0}^{\beta}g=I_{t_0}^{\alpha+\beta}g\).
La misma identidad se extiende a funciones integrables donde las expresiones
están definidas: para $H=T-t_0$, Tonelli da
\[
 \|I^\alpha g\|_{L^1}\leq
 \frac{H^\alpha}{\Gamma(\alpha+1)}\|g\|_{L^1}.
\]
Esta cota permite aplicar Fubini a la composición para $g\in L^1$ y obtener
la igualdad casi en todas partes. En particular, cubre $g=X'$ cuando $X$ es
absolutamente continua. Esta ley para integrales es distinta de
una ley de semigrupo temporal para el estado puntual de un PVI con memoria.

\subsection{Riemann--Liouville y Caputo: diferencia operativa}

Sea \(m-1<q<m\), \(m\in\mathbb N\). Las dos derivadas izquierdas se ordenan de
modo distinto:
\begin{equation}
 {}^{\mathrm{RL}}D_{t_0}^{q}X
 =\frac{\dd^m}{\dd t^m}I_{t_0}^{m-q}X,
 \qquad
 {}^{\mathrm C}D_{t_0}^{q}X
 =I_{t_0}^{m-q}X^{(m)}.
 \label{eq:rl-caputo-general}
\end{equation}
Para \(0<q<1\), la derivada de Caputo de una constante es cero, mientras que
\({}^{\mathrm{RL}}D_{t_0}^{q}c=
c(t-t_0)^{-q}/\Gamma(1-q)\). Caputo admite los datos clásicos
\(X^{(k)}(t_0)\), \(k=0,\ldots,m-1\); Riemann--Liouville conduce en general a
datos fraccionarios. Tampoco se pueden cancelar o componer derivadas
fraccionarias sin revisar los términos de borde. Por eso, en cada modelo se
deben declarar por separado la definición, el terminal inferior, los datos
iniciales y el convenio de memoria
\cite{Podlubny1999,Monje2010}.

Para $0<q<1$, la diferencia entre ambos operadores se expresa exactamente como
\begin{equation}
 {}^{\mathrm{RL}}D_{t_0}^qX
 ={}^{\mathrm C}D_{t_0}^qX+
 \frac{X_0(t-t_0)^{-q}}{\Gamma(1-q)},\qquad t>t_0.
 \label{eq:rl-caputo-initial-bridge}
\end{equation}
Se demuestra separando $X=(X-X_0)+X_0$ y derivando
$I^{1-q}X_0=X_0(t-t_0)^{1-q}/\Gamma(2-q)$.
El mismo PVI puede formularse con Riemann--Liouville si se incorpora ese
término de inicialización. Reemplazar sólo el operador y conservar el
segundo miembro produce, en general, un modelo diferente. Caputo se elige
aquí porque el experimento prescribe valores ordinarios del estado y un
terminal inferior finito. La elección responde a esos datos y al convenio
de memoria, sin establecer una superioridad universal sobre otros operadores.

El comportamiento en el extremo debe distinguirse del dato inicial.
Para $F(X)=v\ne0$, la solución es
\begin{equation}
 X(t)=X_0+\frac{(t-t_0)^q}{\Gamma(q+1)}v,
 \qquad X'(t)=\frac{(t-t_0)^{q-1}}{\Gamma(q)}v.
 \label{eq:caputo-initial-endpoint}
\end{equation}
La derivada ordinaria es integrable y no acotada en el arranque.
La integral beta da ${}^{\mathrm C}D_{t_0}^qX(t)=v$ para todo $t>t_0$;
su límite por la derecha no es cero. Asignar cero a una integral de
intervalo vacío en $t=t_0$ es una convención de extremo y no exige
$F(X_0)=0$. El PVI impone la ecuación para $t>t_0$ en el sentido declarado
y recupera el dato inicial por continuidad de $X$.

\subsection{Derivada de Caputo y forma integral}

Para \(0<q<1\) y \(X:[t_0,T]\to\R^3\) absolutamente continua, la derivada de
Caputo es
\begin{equation}
{}^{\mathrm C}D_{t_0}^{q}X(t)
=\frac{1}{\Gamma(1-q)}
\int_{t_0}^{t}(t-\tau)^{-q}\dot X(\tau)\dd\tau.
\label{eq:caputo}
\end{equation}
La derivada ordinaria se encuentra dentro de la integral y el núcleo
\((t-\tau)^{-q}\) pondera toda la historia entre \(t_0\) y \(t\). Por esta
razón, \(X(t)\) por sí solo no contiene toda la información necesaria para
continuar la misma solución.

Considérese el problema
\begin{equation}
{}^{\mathrm C}D_{t_0}^{q}X(t)=F(t,X(t)),
\qquad X(t_0)=X_0.
\label{eq:pvi_caputo}
\end{equation}
Al aplicar la integral fraccionaria de Riemann--Liouville
\cref{eq:rl-integral}
y usar
\(I_{t_0}^{q}{}^{\mathrm C}D_{t_0}^{q}X
=I_{t_0}^{q}I_{t_0}^{1-q}\dot X
=I_{t_0}^{1}\dot X=X-X_0\), se obtiene la
ecuación de Volterra equivalente
\begin{equation}
X(t)=X_0+\frac{1}{\Gamma(q)}
\int_{t_0}^{t}(t-\tau)^{q-1}F(\tau,X(\tau))\dd\tau.
\label{eq:volterra}
\end{equation}
Recíprocamente, si \(g(t)=F(t,X(t))\) es continua y
\(X-X_0=I_{t_0}^qg\), la ley de composición implica
\[
 \frac{\mathrm d}{\mathrm dt}I_{t_0}^{1-q}(X-X_0)
 =\frac{\mathrm d}{\mathrm dt}I_{t_0}^{1}g=g.
\]
Esta expresión define la derivada de Caputo en la clase integral de
soluciones continuas \(X-X_0\in I_{t_0}^q(C)\); coincide con
\eqref{eq:caputo} cuando \(X\) es absolutamente continua.
La formulación integral permite así tratar el PVI sin atribuir regularidad
clásica adicional a toda integral fraccionaria de una función continua.
Para campos medibles en el tiempo, como una continuación por etapas,
se usa la clase $X_0+I_{t_0}^q(L^\infty)$ sobre intervalos compactos y
${}^{\mathrm C}D^qX=(\mathrm d/\mathrm dt)I^{1-q}(X-X_0)$ casi en todas
partes. En esta clase $I^{1-q}(X-X_0)=I^1g$ es absolutamente continua.
Recíprocamente, integrar la derivada generalizada da esa identidad, con
constante nula en $t_0$. Si $h=X-X_0-I^qg$, entonces $I^{1-q}h=0$;
aplicar $I^q$ y derivar $I^1h=0$ da $h=0$ casi en todas partes y, por
continuidad, en todo el intervalo. Esta prueba establece ambas direcciones
de la equivalencia sin suponer la existencia de $X'$ en toda la clase integral
\cite{Gomoyunov2018Convex}.
La ecuación \eqref{eq:volterra} es la base del integrador ABM y de la continuación con memoria:
en ambos casos se discretiza una sola integral cuyo límite inferior permanece
fijo.

La existencia local se obtiene directamente por contracción. En una bola
cerrada $\overline B(X_0,R)$ contenida en el dominio, supóngase
$\|F(t,X)\|\le M$ y Lipschitz espacial uniforme con constante $L$.
El operador $\mathcal TX=X_0+I^q[F(\cdot,X)]$ satisface
\[
 \|\mathcal TX-X_0\|_\infty\le\frac{M\tau^q}{\Gamma(q+1)},
 \qquad
 \|\mathcal TX-\mathcal TY\|_\infty
 \le\frac{L\tau^q}{\Gamma(q+1)}\|X-Y\|_\infty.
\]
Elegir el primer coeficiente menor o igual que $R$ y el segundo estrictamente
menor que uno permite aplicar el teorema de Banach sobre las curvas continuas
que permanecen en la bola. Para tiempos posteriores se conserva la historia
acumulada. Si una solución permanece en un compacto del dominio hasta un
tiempo máximo finito, la acotación del campo y la continuidad de su integral
fraccionaria permiten extenderla con esa historia; por tanto, no puede
terminar allí su intervalo maximal. Este argumento fundamenta el uso de
cotas absorbentes para demostrar existencia global.

La verificación del esquema ABM de \cref{sec:abm} incluye resolver una ecuación
lineal cuya solución se exprese mediante
\cref{eq:mittag-leffler-definition}, derivar los pesos del predictor y del
corrector y comparar \(h\), \(h/2\) y \(h/4\). Una aproximación de memoria
corta requiere identificar la parte eliminada de la integral y acotar el error;
por tanto, modifica el modelo numérico además de su implementación.

\subsection{Transformada de Laplace de la derivada de Caputo}

Supóngase que \(X\) es absolutamente continua en intervalos compactos y
que \(X\) y \(\dot X\) tienen crecimiento a lo sumo exponencial, de modo
que sus transformadas y la convolución convergen en un semiplano derecho.
La relación en Laplace se obtiene de \cref{eq:caputo}. Para hacer
explícito el instante de inicio, se introduce
\[
u=t-t_0,\qquad \widetilde X(u)=X(t_0+u),\qquad u\geq0.
\]
Entonces
\[
{}^{\mathrm C}D_{0}^{q}\widetilde X
=I_0^{1-q}\widetilde X',
\]
y la propiedad de convolución da
\begin{align}
\mathcal L_u\!\left\{I_0^{1-q}\widetilde X'\right\}(s)
&=s^{-(1-q)}
\mathcal L_u\!\left\{\widetilde X'\right\}(s)\nonumber\\
&=s^{q-1}\left[s\widehat{\widetilde X}(s)-\widetilde X(0)\right].
\end{align}
Por tanto,
\begin{equation}
\boxed{
\mathcal L_u\!\left\{{}^{\mathrm C}D_0^q\widetilde X(u)\right\}(s)
=s^q\widehat{\widetilde X}(s)-s^{q-1}X(t_0).
}
\label{eq:laplace_caputo}
\end{equation}
El término \(s^{q-1}X(t_0)\) contiene la condición inicial. La función de
transferencia resulta de aplicar
\cref{eq:laplace_caputo} al bloque lineal y separar la respuesta debida a la
entrada.

La misma identidad de Laplace vale en la clase generalizada si
$g={}^{\mathrm C}D^qX$ tiene integrabilidad exponencial y
$X-X_0=I^qg$. Basta que $e^{-\sigma t}g\in L^1([0,\infty))$ para algún
$\sigma>0$: Fubini con ese peso justifica la convolución, y tanto
$X-X_0$ como $H=I^{1-q}(X-X_0)=I^1g$ y $H'=g$ tienen las
transformadas necesarias en un semiplano derecho. En efecto,
\[
 \mathcal L\!\left\{\frac{\mathrm d}{\mathrm dt}
 I^{1-q}(X-X_0)\right\}
 =s\,s^{q-1}\left(\widehat X-\frac{X_0}{s}\right),
\]
porque $I^{1-q}(X-X_0)$ se anula en el extremo.
La propiedad de convolución exige integrabilidad absoluta con el peso
exponencial; no se deduce de una existencia global que aún no esté probada.
La rama es $s^q=\exp(q\operatorname{Log}s)$, con
$\operatorname{Arg}s\in(-\pi,\pi)$.

\subsection{Liouville--Weyl como predictor armónico}

Un sistema autónomo de Caputo con límite inferior fijo no posee, en general,
soluciones periódicas exactas no constantes
\cite{TavazoeiHaeri2009,AreaLosadaNieto2014}. Sin embargo, la construcción de
una semilla requiere evaluar la respuesta a un armónico. Para ello se utiliza
un problema auxiliar estacionario definido sobre una historia extendida hacia
\(-\infty\). En el sentido de Fourier, el operador de Liouville--Weyl satisface
\cite{Vigue2019}
\begin{equation}
{}^{\mathrm{LW}}D_t^q\e^{\ii\omega t}
=(\ii\omega)^q\e^{\ii\omega t},
\qquad
(\ii\omega)^q
=\omega^q\e^{\ii q\pi/2},
\quad \omega>0.
\label{eq:lw_armonico}
\end{equation}
Se adopta la rama principal:
\begin{equation}
(\ii\omega)^q
=\omega^q\left[
\cos\left(\frac{q\pi}{2}\right)
+\ii\sin\left(\frac{q\pi}{2}\right)
\right].
\label{eq:rama_principal}
\end{equation}

La relación entre el problema auxiliar y la dinámica causal puede escribirse
de forma explícita. Sea \(\overline X(t)=X_-(t)\) para \(t\leq t_0\) y
\(\overline X(t)=X(t)\) para \(t\geq t_0\), con continuidad en \(t_0\).
Supóngase además que ambas ramas son localmente absolutamente continuas
y que converge la integral impropia que aparece a continuación.
Para una historia armónica se utiliza la convergencia oscilatoria o su
regularización de Abel, precisada en el desarrollo del balance.
En este dominio,
\begin{equation}
\begin{aligned}
{}_{-\infty}^{\mathrm{CW}}D_t^q\overline X(t)
&=
\frac{1}{\Gamma(1-q)}
\int_{-\infty}^{t}(t-\tau)^{-q}\dot{\overline X}(\tau)\dd\tau\\
&={}^{\mathrm C}D_{t_0}^{q}X(t)+\mathcal E_{X_-}(t),
\end{aligned}
\label{eq:puente_weyl_caputo}
\end{equation}
donde
\begin{equation}
\mathcal E_{X_-}(t)=
\frac{1}{\Gamma(1-q)}
\int_{-\infty}^{t_0}(t-\tau)^{-q}\dot X_-(\tau)\dd\tau.
\label{eq:memoria_previa}
\end{equation}
Si \(\dot X_-\in L^1(-\infty,t_0)\),
\begin{equation}
\|\mathcal E_{X_-}(t)\|
\leq
\frac{\|\dot X_-\|_{L^1(-\infty,t_0)}}{\Gamma(1-q)}
(t-t_0)^{-q}.
\label{eq:cota_memoria_previa}
\end{equation}
La ecuación armónica de Liouville--Weyl proporciona así el predictor
frecuencial. La trayectoria que se evalúa y sondea se calcula después con
Caputo y un límite inferior finito.

\subsection{Recuperación del orden entero a tiempo finito}
\label{sec:limite-entero-tiempo-finito}

Las identidades algebraicas del campo, como $F(e)=0$, no dependen del
orden. Para comparar las trayectorias sí se necesita una estimación.
Sean $X_q$ y $X_1$ soluciones con el mismo dato inicial en $[0,T]$,
contenidas en un compacto común donde $\|F\|\le M$ y $F$ sea Lipschitz
con constante $L$. Defina
\[
 k_q(t)=\frac{t^{q-1}}{\Gamma(q)},\qquad
 \varepsilon_q=\int_0^T|k_q(t)-1|\,\mathrm dt.
\]
Si $q\in[q_*,1]$, $q_*>0$, el integrando está dominado por una constante
por $t^{q_*-1}+1$, integrable en $[0,T]$. La convergencia puntual
$k_q(t)\to1$ para $t>0$ implica $\varepsilon_q\to0$ por convergencia
dominada. Restar las ecuaciones de Volterra y usar la cota Lipschitz da
\begin{align*}
 \|X_q(t)-X_1(t)\|
 &\le M\varepsilon_q+
 L\int_0^t k_q(t-s)\|X_q(s)-X_1(s)\|\,\mathrm ds,\\
 \sup_{0\le t\le T}\|X_q(t)-X_1(t)\|
 &\le M\varepsilon_q E_q(LT^q)\longrightarrow0.
\end{align*}
La segunda desigualdad se obtiene iterando la primera: la composición
de integrales produce el término $(Lt^q)^j/\Gamma(qj+1)$.
El factor $E_q(LT^q)$ permanece uniformemente acotado. Para comprobarlo,
ponga $B=L\max(T,T^{q_*})$; la cola de su serie está dominada por
$\sum_j B^j/\Gamma(q_*j+1)$, convergente, porque $\Gamma$ es creciente
para argumentos suficientemente grandes. Los términos restantes son
finitos y continuos en $q$.

La hipótesis de compacto común se comprueba mediante una cota uniforme o
se usa el argumento hasta un primer tiempo de salida. El resultado
justifica una comparación sobre una ventana fija; no permite intercambiar
los límites $q\to1$ y $T\to\infty$. La persistencia de atractores,
cuencas u ocultedad exige control asintótico adicional. Asimismo, la
continuación de una raíz del balance armónico requiere que su Jacobiano
respecto de las incógnitas sea invertible, y la continuación causal
conserva la historia. Son requisitos distintos de la continuidad a
tiempo finito demostrada aquí.
\fi

\ifdefined\HAFOFractionalBlockStability
\ifdefined\HAFOFractionalBlockOperators\else
\ifHAFOBasinMaterial
\section{Estabilidad fraccionaria y cuencas con memoria}
\else
\section{Estabilidad fraccionaria}
\fi
\label{sec:fundamentos-estabilidad}
\fi
\subsection{Linealización y criterio de Matignon}

Sea \(e\) un equilibrio de
\[
{}^{\mathrm C}D_{t_0}^{q}X=F(X),\qquad F(e)=0.
\]
Con \(X=e+\xi\) y \(F\) diferenciable en \(e\),
\[
F(e+\xi)=J_e\xi+o(\|\xi\|),
\qquad
J_e=\left.\frac{\partial F}{\partial X}\right|_{X=e}.
\]
El resto es cuadrático si se supone, por ejemplo, \(F\in C^{1,1}\) localmente
o \(F\in C^2\). La diferenciabilidad por sí sola no basta: para
\(F(x)=|x|^{3/2}\), \(F'(0)=0\), pero el resto no es
\(\mathcal O(|x|^2)\). En un modelo PWL, el Jacobiano clásico se usa cuando el
equilibrio pertenece al interior de una región afín; un equilibrio sobre una
superficie de conmutación requiere un análisis específico.
La aproximación lineal es
\begin{equation}
{}^{\mathrm C}D_{t_0}^{q}\xi=J_e\xi.
\label{eq:lineal_caputo}
\end{equation}
Para orden conmensurable \(0<q<1\), el sistema lineal es asintóticamente
estable si y sólo si todos los valores propios de \(J_e\) son no nulos y
cumplen el criterio de Matignon \cite{Matignon1996}
\begin{equation}
|\arg\lambda_j(J_e)|>\frac{q\pi}{2}.
\label{eq:matignon}
\end{equation}
La solución del problema lineal es
\[
 \xi(t)=E_q\!\left(J_e(t-t_0)^q\right)\xi(t_0),
 \qquad
 E_q(J_e\tau^q)=\sum_{k=0}^{\infty}
 \frac{J_e^k\tau^{qk}}{\Gamma(qk+1)}.
\]
Esta expresión se obtiene de
\(\widehat\xi(s)=(s^qI-J_e)^{-1}s^{q-1}\xi(t_0)\), al trasladar
\(t_0\) a cero. El criterio de Matignon caracteriza el decaimiento de
esta solución, incluidos los bloques de Jordan; por eso se aplica al
espectro completo, no sólo a la traza o al determinante de \(J_e\).
Para trasladar esta conclusión al equilibrio no lineal se aplica el teorema de
linealización correspondiente. Una hipótesis suficiente es \(F\in C^1\) en
una vecindad de \(e\). Para comprobar la condición sobre el resto, defina
\(R(\xi)=F(e+\xi)-J_e\xi\) y escoja una bola convexa de radio
\(\delta\) contenida en esa vecindad. El teorema fundamental del cálculo da
\begin{align*}
 R(\xi)-R(\eta)
 &=\int_0^1\bigl[DF(e+\eta+s(\xi-\eta))-J_e\bigr]
                (\xi-\eta)\,\mathrm ds,\\
 \|R(\xi)-R(\eta)\|
 &\leq\sup_{\|u\|\leq\delta}\|DF(e+u)-J_e\|\,
          \|\xi-\eta\|.
\end{align*}
La continuidad de \(DF\) hace que el supremo tienda a cero cuando
\(\delta\to0\). Ésta es la pequeñez Lipschitz requerida por el teorema
no lineal. Bajo el criterio angular
estricto, el equilibrio es localmente estable en el sentido de
Mittag--Leffler \cite{CongTuanTrinh2020Asymptotic}.
El mecanismo de la estabilidad asintótica puede explicitarse mediante
variación de constantes \cite{CongDoanSiegmundTuan2016}. Bajo el sector
estricto, las estimaciones de Mittag--Leffler proporcionan
\[
 M_J=\sup_{t\ge0}\|E_q(J_et^q)\|<\infty,\qquad
 C_J=\int_0^\infty\|t^{q-1}E_{q,q}(J_et^q)\|\,\mathrm dt<\infty.
\]
El núcleo es $O(t^{q-1})$ en cero y $O(t^{-q-1})$ en el infinito.
Sea $\ell_\delta$ la constante Lipschitz del resto en la bola anterior.
Se elige $\delta$ con $C_J\ell_\delta<1$ y un dato inicial tal que
$M_J\|\xi_0\|<(1-C_J\ell_\delta)\delta$. La representación
\[
 \xi(t)=E_q(J_et^q)\xi_0+
 \int_0^t(t-s)^{q-1}E_{q,q}(J_e(t-s)^q)R(\xi(s))\,\mathrm ds
\]
excluye un primer alcance de la frontera y da estabilidad local.
Para la atracción, se divide la convolución en un intervalo inicial fijo,
cuyo aporte tiende a cero, y la cola. Si
$a=\limsup_{t\to\infty}\|\xi(t)\|$, la integrabilidad del núcleo da
$a\le C_J\ell_\delta a$; por tanto, $a=0$.
La fórmula integral puede verificarse en horizontes finitos mediante la
serie de Volterra: los términos están dominados por
$\|J_e\|^jT^{qj}/\Gamma(qj+1)$ y, para la fuerza acotada, por
$\|J_e\|^jT^{q(j+1)}/\Gamma(q(j+1)+1)$. La convergencia uniforme
permite integrar término a término y la unicidad identifica la solución.
Se define el margen angular
\begin{equation}
m_M(e)=
\min_j\left(
|\arg\lambda_j(J_e)|-\frac{q\pi}{2}
\right).
\label{eq:margen_matignon}
\end{equation}
Cuando todos los autovalores son no nulos, un margen positivo satisface el
criterio estricto de estabilidad lineal y un margen negativo identifica al
menos un modo lineal estrictamente inestable. El caso \(m_M=0\) exige estudiar
la frontera angular y los términos no lineales. Si existe un autovalor cero,
su argumento no está definido y el criterio angular no decide su contribución
a la estabilidad.

\ifHAFOBasinMaterial
\subsection{Cuenca con convenio común de memoria}

Para comparar condiciones iniciales en Caputo se fija el mismo \(t_0\) y se
inicia cada problema con la historia definida por su propio valor \(X_0\).
Con este convenio de memoria, la cuenca se define por
\begin{equation}
\basin(\mathcal A)=
\left\{
X_0\in\R^3:
\operatorname{dist}\bigl(X(t;t_0,X_0),\mathcal A\bigr)\to0
\ \text{cuando}\ t\to\infty
\right\}.
\label{eq:cuenca_reset}
\end{equation}
La fijación del terminal inferior y del perfil inicial permite comparar los
destinos de los reinicios alrededor de los equilibrios. La convergencia
proyectada de \cref{eq:cuenca_reset} debe distinguirse de la convergencia
de perfiles en un espacio funcional.
\fi
\fi
\endgroup
\fi
